\documentclass[a4paper,12pt,oneside]{book}

\usepackage{fontspec}
\setmainfont[
    BoldFont = Arial Bold,
    ItalicFont = Arial Italic,
    BoldItalicFont = Arial Bold Italic
]{Arial}
\renewcommand{\familydefault}{\sfdefault}


%--- Encoding & Language --------------------------------------------------
\usepackage[T1]{fontenc}
\usepackage[utf8]{inputenc}
\usepackage[italian]{babel}

%--- Typography -----------------------------------------------------------
\usepackage{lmodern}
\usepackage{microtype}
%--- Immagini-----------------------------------------------------------
\usepackage{graphicx}
\usepackage{float}
%--- Page Layout ----------------------------------------------------------
\usepackage[a4paper, top=2.5cm, bottom=2.8cm, left=3cm, right=2.5cm,
            headheight=15pt]{geometry}

%--- Colors (deve essere caricato prima di fancyhdr) ----------------------
\usepackage{xcolor}
\definecolor{cvtblue}{RGB}{14,73,143}
\definecolor{cvtlightblue}{RGB}{228,238,255}
\definecolor{cvtgray}{RGB}{80,80,80}
\definecolor{cvtgreen}{RGB}{0,120,50}
\definecolor{cvtred}{RGB}{160,30,30}
\definecolor{cvtamber}{RGB}{180,120,0}
\definecolor{codebg}{RGB}{246,248,250}
\definecolor{codeline}{RGB}{14,73,143}
\definecolor{keycolor}{RGB}{0,0,180}
\definecolor{sectioncolor}{RGB}{140,30,30}
\definecolor{commentcolor}{RGB}{110,110,110}

%--- Headers & Footers ----------------------------------------------------
\usepackage{fancyhdr}
\pagestyle{fancy}
\fancyhf{}
\fancyhead[L]{\footnotesize\textcolor{cvtgray}{\nouppercase{\leftmark}}}
\fancyhead[R]{\footnotesize\textcolor{cvtgray}{CVTScript v3}}
\fancyfoot[C]{\small\thepage}
\renewcommand{\headrulewidth}{0.4pt}

%--- Listings -------------------------------------------------------------
\usepackage{listings}

\lstdefinelanguage{CVTScript}{
  sensitive=true,
  morecomment=[l]{\#},
  moredelim=[s][\color{sectioncolor}\bfseries]{[}{]},
}

\lstdefinestyle{cvt}{
  language=CVTScript,
  basicstyle=\ttfamily\small,
  backgroundcolor=\color{codebg},
  frame=l,
  framerule=3pt,
  rulecolor=\color{codeline},
  commentstyle=\color{commentcolor}\itshape,
  breaklines=true,
  keepspaces=true,
  columns=flexible,
  xleftmargin=10pt,
  xrightmargin=4pt,
  aboveskip=6pt,
  belowskip=6pt,
  showstringspaces=false,
  extendedchars=true,
  literate=
    {à}{{\`a}}1 {è}{{\`e}}1 {ì}{{\`i}}1 {ò}{{\`o}}1 {ù}{{\`u}}1
    {À}{{\`A}}1 {È}{{\'E}}1 {Ì}{{\`I}}1 {Ò}{{\`O}}1 {Ù}{{\`U}}1
    {é}{{\'e}}1 {á}{{\'a}}1 {ó}{{\'o}}1 {ú}{{\'u}}1
    {ë}{{\"e}}1 {ï}{{\"i}}1 {ö}{{\"o}}1 {ü}{{\"u}}1,
}
\lstset{style=cvt}

%--- Boxes ----------------------------------------------------------------
\usepackage[most,breakable]{tcolorbox}

\newtcolorbox{esempiocvt}[1][Esempio]{
  enhanced, breakable,
  colback=codebg, colframe=cvtblue,
  coltitle=white, fonttitle=\bfseries\small,
  title={#1}, arc=3pt, boxrule=0.8pt,
  left=2pt, right=2pt, top=0pt, bottom=2pt,
  toptitle=3pt, bottomtitle=3pt,
}

\newtcolorbox{notablu}[1][Nota]{
  enhanced, breakable,
  colback=cvtlightblue, colframe=cvtblue!50,
  fonttitle=\bfseries\small, title={#1},
  arc=3pt, boxrule=0.6pt,
  left=6pt, right=6pt, top=4pt, bottom=4pt,
}

\newtcolorbox{attenzione}[1][Attenzione]{
  enhanced, breakable,
  colback=orange!8, colframe=cvtamber,
  fonttitle=\bfseries\small\color{cvtamber}, title={#1},
  arc=3pt, boxrule=0.8pt,
  left=6pt, right=6pt, top=4pt, bottom=4pt,
}

\newtcolorbox{regola}[1][Regola d'oro]{
  enhanced, breakable,
  colback=green!5, colframe=cvtgreen,
  fonttitle=\bfseries\small\color{cvtgreen}, title={#1},
  arc=3pt, boxrule=0.8pt,
  left=6pt, right=6pt, top=4pt, bottom=4pt,
}

\newtcolorbox{trucchetto}[1][Trucchetto]{
  enhanced, breakable,
  colback=yellow!10, colframe=cvtamber!70,
  fonttitle=\bfseries\small\color{cvtamber}, title={#1},
  arc=3pt, boxrule=0.6pt,
  left=6pt, right=6pt, top=4pt, bottom=4pt,
}

\newtcolorbox{sbagliato}[1][Sbagliato]{
  enhanced, breakable,
  colback=red!5, colframe=cvtred,
  fonttitle=\bfseries\small\color{cvtred}, title={#1},
  arc=3pt, boxrule=0.8pt,
  left=6pt, right=6pt, top=4pt, bottom=4pt,
}

\newtcolorbox{giusto}[1][Giusto]{
  enhanced, breakable,
  colback=green!5, colframe=cvtgreen,
  fonttitle=\bfseries\small\color{cvtgreen}, title={#1},
  arc=3pt, boxrule=0.8pt,
  left=6pt, right=6pt, top=4pt, bottom=4pt,
}

%--- Tables ---------------------------------------------------------------
\usepackage{booktabs}
\usepackage{array}
\usepackage{tabularx}
\usepackage{longtable}

%--- Misc -----------------------------------------------------------------
\usepackage{amssymb}
\usepackage{enumitem}
\usepackage{parskip}
\usepackage{float}
\usepackage{caption}

%--- Chapter & Section style ----------------------------------------------
\usepackage{titlesec}
\titleformat{\chapter}[display]
  {\normalfont\huge\bfseries\color{cvtblue}}
  {\chaptertitlename\ \thechapter}{10pt}{\Huge}
\titleformat{\section}
  {\normalfont\Large\bfseries\color{cvtblue}}{\thesection}{1em}{}
\titleformat{\subsection}
  {\normalfont\large\bfseries\color{cvtgray}}{\thesubsection}{1em}{}

%--- Hyperlinks -----------------------------------------------------------
\usepackage[colorlinks=true, linkcolor=cvtblue,
            urlcolor=cvtblue, citecolor=cvtblue]{hyperref}

%--- Utility commands -----------------------------------------------------
\newcommand{\kw}[1]{\texttt{\color{keycolor}#1}}
\newcommand{\val}[1]{\texttt{\color{cvtgreen}#1}}
\newcommand{\sez}[1]{\texttt{\color{sectioncolor}\textbf{#1}}}
\newcommand{\file}[1]{\texttt{#1}}
\newcommand{\siok}{\textcolor{cvtgreen}{\textbf{\checkmark}}}
\newcommand{\sino}{\textcolor{cvtred}{\textbf{$\times$}}}

%--- Placeholder per immagini da inserire in seguito ----------------------
% Uso: \imgph{larghezza}{descrizione del soggetto da disegnare}
% Quando l'immagine sarà pronta, sostituire la chiamata con:
%   \includegraphics[width=...]{immagini/capN/nomefile.png}
\newcommand{\imgph}[2]{%
  \begin{tcolorbox}[colback=cvtlightblue!30,colframe=cvtblue!40,%
    width=#1,halign=center,valign=center,height=5cm,boxrule=0.5pt,arc=3pt]%
    \textcolor{cvtgray}{\itshape\footnotesize\textbf{[Immagine da inserire]}\\[3pt]#2}%
  \end{tcolorbox}}

%==========================================================================
\begin{document}
%==========================================================================

\frontmatter
\pagestyle{plain}

%--- Cover page ------------------------------------------------------------
\begin{titlepage}
  \pagecolor{cvtblue}\color{white}
  \vspace*{1.5cm}
  \begin{center}
    {\large\textsf{\textbf{CAMPUS VIRTUAL TRAINING}}}\\[0.4cm]
    \textcolor{white!60!cvtblue}{\rule{\textwidth}{0.6pt}}\\[2cm]
    {\fontsize{30}{36}\selectfont\textbf{CVTScript v3}}\\[0.5cm]
    {\fontsize{18}{22}\selectfont Manuale per Principianti}\\[0.3cm]
    {\fontsize{13}{16}\selectfont Come scrivere un tutorial passo-passo,\\spiegato semplice}\\[2.5cm]
    \textcolor{white!60!cvtblue}{\rule{\textwidth}{0.4pt}}\\[1cm]
    \begin{tabular}{@{}ll@{}}
      \textbf{Versione del linguaggio:} & 3.0 (completa) \\[5pt]
      \textbf{Data:}                    & Maggio 2026 \\[5pt]
      \textbf{A chi serve:}             & A chi non sa programmare \\[5pt]
      \textbf{Tempo di lettura:}        & 50--60 minuti
    \end{tabular}
  \end{center}
  \vfill
  \begin{center}
    {\footnotesize Leggi una volta, prova subito, divertiti.}
  \end{center}
\end{titlepage}
\pagecolor{white}\color{black}

\chapter*{Per cominciare}

Questo manuale ti insegna a scrivere un \textbf{tutorial 3D} per Campus Virtual Training
\textbf{senza saper programmare}. Lo leggi una volta, ti diverti, e alla fine sei
capace di creare le tue lezioni interattive.

\medskip
\textbf{Cosa NON serve.} Non serve conoscere nessun linguaggio di programmazione.
Non serve saper usare strumenti tecnici. Non serve avere studiato informatica.

\medskip
\textbf{Cosa serve invece.} Serve sapere \emph{cosa} vuoi insegnare allo studente,
passo dopo passo. Esattamente come quando spieghi a un amico come fare una torta:
prima si prende la ciotola, poi si rompono le uova, poi si aggiunge la farina\ldots

\medskip
\textbf{Come è organizzato il manuale.} Sono \textbf{capitoli brevi}, divisi in
tre parti. La \textbf{prima parte} (cap.\ 1--10) spiega le \textbf{basi}: ti
basta per scrivere un tutorial completo. La \textbf{seconda parte}
(cap.\ 11--17) spiega le \textbf{cose avanzate} (oggetti in mano, finestre con
messaggi, animazioni, oggetti interattivi). La \textbf{terza parte}
(cap.\ 18) spiega i \textbf{file di configurazione} che vivono attorno al
tutorial. Alla fine c'è un \textbf{glossario} con tutte le parole importanti
tradotte in italiano semplice. Se ti perdi, torna al glossario.

\begin{figure}[H]
  \centering
  % \includegraphics[width=12cm]{immagini/frontmatter/mappa_manuale.png}
  \imgph{12cm}{Mappa visiva del manuale: tre blocchi orizzontali colorati ``Basi 1--10'', ``Avanzato 11--17'', ``Configurazione 18'', con frecce di lettura e icone tematiche.}
  \caption{Le tre parti del manuale: leggile in ordine la prima volta, poi consulta a salti.}
  \label{fig:mappa-manuale}
\end{figure}

\vfill
\begin{center}
\textit{Buona lettura!}
\end{center}
\clearpage

\tableofcontents

%==========================================================================
\mainmatter
\pagestyle{fancy}
%==========================================================================

%--------------------------------------------------------------------------
\chapter{Cos'è un tutorial?}
\label{ch:intro}
%--------------------------------------------------------------------------

Pensa a quando insegni a un amico \textbf{come fare una torta}. Non gli dici tutto
in una volta. Vai per piccoli passi:

\begin{itemize}
  \item \textbf{Passo 1}: prendi una ciotola.
  \item \textbf{Passo 2}: rompi due uova dentro.
  \item \textbf{Passo 3}: aggiungi la farina.
  \item \ldots
\end{itemize}

Un tutorial nel programma Campus Virtual Training funziona \textbf{esattamente così},
ma invece della torta lo studente sta imparando a \textbf{smontare una pompa},
\textbf{cambiare un filtro} o \textbf{lubrificare un elettromandrino}.

\begin{figure}[H]
  \centering
  % \includegraphics[width=13cm]{immagini/cap01/torta_vs_pompa.png}
  \imgph{13cm}{Confronto a due colonne: a sinistra una sequenza ``ciotola $\rightarrow$ uova $\rightarrow$ farina $\rightarrow$ forno'' (la torta), a destra ``vite $\rightarrow$ coperchio $\rightarrow$ filtro $\rightarrow$ rimontaggio'' (la pompa). Sopra: la scritta ``stessa logica, mondi diversi''.}
  \caption{Stessa logica, due mondi diversi: la metafora della ricetta funziona perché ogni tutorial è una sequenza di passi semplici.}
  \label{fig:torta-vs-pompa}
\end{figure}

Tu, nel file \file{tutorial.cvtscript}, scrivi la lista dei passi. Il programma
li esegue uno alla volta e fa muovere i pezzi sullo schermo.

\begin{figure}[H]
  \centering
  % \includegraphics[width=12cm]{immagini/cap01/interfaccia_annotata.png}
  \imgph{12cm}{Screenshot del tutorial in esecuzione con quattro callout numerati: 1=modello 3D al centro, 2=fumetto della spiegazione, 3=barra strumenti in basso, 4=freccia ``avanti'' a destra.}
  \caption{L'interfaccia che vede lo studente: i quattro elementi che si ripetono in ogni step.}
  \label{fig:interfaccia-annotata}
\end{figure}

\begin{notablu}[In poche parole]
Tu sei l'\textbf{autore di una ricetta}. Lo studente è un \textbf{cuoco apprendista}.
Il programma è il \textbf{cuoco esperto} che dimostra ogni passaggio. Tu scrivi la
ricetta in modo chiaro, e il programma ci pensa a far vedere tutto.
\end{notablu}

%--------------------------------------------------------------------------
\chapter{I modelli 3D: la scena dove tutto accade}
\label{ch:modelli3d}
%--------------------------------------------------------------------------

Adesso che sai cos'è un tutorial, c'è una cosa da capire prima di
iniziare a scrivere il tuo primo file: i pezzi che lo studente
\textbf{vede e manipola} sullo schermo sono \textbf{modelli 3D}, ovvero
oggetti tridimensionali caricati dal programma. Non li crei tu scrivendo
il tutorial --- esistono già come file separati. Tu nel tutorial dici
solo \emph{``prendi quella vite, svitala, mettila giù''} e il programma
sa già come è fatta la vite perché l'ha caricata all'avvio.

\begin{notablu}[Puoi saltare questo capitolo?]
Se qualcun altro ha già preparato i modelli 3D e sai solo che si trovano
nella cartella \file{models/}, puoi passare direttamente al
Capitolo~\ref{ch:struttura}. Torna qui quando vuoi capire come vengono
creati o se devi aggiungerli tu.
\end{notablu}

\section{Il formato richiesto: \texttt{.glb}}

Campus Virtual Training usa \textbf{Three.js} come motore di rendering
e carica esclusivamente modelli in formato \textbf{glTF Binary}
(\file{.glb}). È uno standard aperto e offre tre vantaggi pratici:

\begin{itemize}
  \item \textbf{File singolo}: geometria, materiali e texture sono
        contenuti in un unico file binario.
  \item \textbf{Compatto}: la codifica binaria riduce le dimensioni
        rispetto a formati testuali come OBJ o COLLADA.
  \item \textbf{Compatibilità web}: è il formato nativo di Three.js e
        non richiede conversioni a runtime.
\end{itemize}

\begin{notablu}[Ricorda]
Se hai modelli in altri formati (\file{.step}, \file{.stl}, \file{.obj},
\file{.fbx}), devi convertirli in \file{.glb} prima di poterli usare.
\textbf{Blender} è lo strumento consigliato per la conversione.
\end{notablu}

\begin{figure}[H]
  \centering
  % \includegraphics[width=12cm]{immagini/cap02/formati_confronto.png}
  \imgph{12cm}{Confronto visivo formati 3D: quattro icone-file in fila (STEP, OBJ, FBX, GLB) con badge di idoneità. Solo GLB ha la spunta verde; sotto ogni formato tre micro-icone (file singolo / compatto / web-ready) accese o spente.}
  \caption{Perché GLB: un solo file, dimensioni ridotte, nessuna conversione a runtime.}
  \label{fig:formati-confronto}
\end{figure}

\section{Esportazione da Blender}

Blender è il software consigliato per preparare i modelli 3D. La
procedura tipica è:

\begin{enumerate}
  \item \textbf{Importa} il modello CAD originale (STEP, IGES o STL sono
        i più comuni nelle officine meccaniche).
  \item \textbf{Pulisci la mesh}: elimina facce duplicate, correggi le
        normali invertite e riduci la complessità poligonale dove
        possibile.
  \item \textbf{Assegna i materiali}: ogni materiale definito in Blender
        verrà esportato nel GLB. Usa nomi significativi.
  \item \textbf{Posiziona l'origine}: l'origine dell'oggetto in Blender
        corrisponde al \textbf{pivot} attorno al quale Campus Virtual
        Training calcolerà rotazioni e traslazioni.
  \item \textbf{Esporta} con \emph{File $\rightarrow$ Export $\rightarrow$
        glTF 2.0 (.glb/.gltf)}, selezionando il formato \file{.glb}.
        Nel pannello di esportazione, sotto la sezione \emph{Geometry},
        abilita l'opzione \emph{Compression} (compressione \textbf{Draco}).
        Riduce le dimensioni dei file fino al 60--80\,\% --- molto utile
        per scenari con molti modelli come la Pompa Becker (oltre 40 file
        GLB). Impostazioni consigliate:
    \begin{itemize}
      \item \emph{Compression Level}: 6 --- buon compromesso tra
            dimensione e velocità di decodifica.
      \item \emph{Quantization Position}: 14 bit --- precisione sufficiente
            per componenti meccanici.
      \item \emph{Quantization Normal}: 10 bit.
    \end{itemize}
\end{enumerate}

\begin{trucchetto}
Per i materiali \textbf{emissivi} (LED, spie luminose), usa nomi che
contengano \texttt{luminoso}, \texttt{emissive} oppure \texttt{glow}:
il sistema li rileverà automaticamente e ne potenzierà la luminosità
(\texttt{emissiveIntensity = 3.0}).
\end{trucchetto}

\begin{figure}[H]
  \centering
  % \includegraphics[width=12cm]{immagini/cap02/blender_export.png}
  \imgph{12cm}{Screenshot del pannello di esportazione di Blender (File $\rightarrow$ Export $\rightarrow$ glTF 2.0). Sezione \emph{Geometry} aperta con la voce \emph{Compression} cerchiata in rosso e i tre slider Position/Normal/UV visibili.}
  \caption{Il pannello di esportazione di Blender. La compressione Draco si trova sotto Geometry $\rightarrow$ Compression.}
  \label{fig:blender-export}
\end{figure}

\subsection{Convenzioni di denominazione}

La denominazione dei file è fondamentale. Ogni file GLB deve avere un
nome \textbf{univoco}, \textbf{tutto in minuscolo} e \textbf{senza spazi}.
Per gli scenari della Pompa Becker, ad esempio:

\begin{itemize}
  \item \textbf{Componente principale}: \file{corpo.glb} --- il corpo
        della pompa, elemento statico di riferimento.
  \item \textbf{Coperchio}: \file{coperchio.glb} --- il coperchio che
        verrà smontato e rimontato durante il tutorial.
  \item \textbf{Parti interne}: \file{filtro.glb}, \file{culatta.glb},
        \file{flangia.glb} --- componenti accessibili dopo lo smontaggio.
  \item \textbf{Viteria}: \file{vite\_coperchio\_1.glb},
        \file{vite\_coperchio\_2.glb}\ldots{} --- ogni vite è un file
        separato per poterla animare indipendentemente.
  \item \textbf{Ambiente}: \file{pavimento.glb} --- base statica della
        scena per dare contesto visivo.
\end{itemize}

\begin{attenzione}
Ogni oggetto che lo studente deve manipolare direttamente (svitare,
estrarre, trascinare) deve essere esportato come file \file{.glb}
autonomo. I sotto-assiemi che si muovono \emph{automaticamente} (per
esempio gli assi di una macchina) possono invece restare come figli
all'interno del GLB principale e venire animati tramite la notazione
puntata \kw{element = a500.porta}.
\end{attenzione}

\begin{figure}[H]
  \centering
  % \includegraphics[width=11cm]{immagini/cap02/cartella_models.png}
  \imgph{11cm}{Screenshot dell'Esplora risorse (o Finder) aperto sulla cartella \file{models/} della Pompa Becker. Si vedono i file GLB con nomi minuscoli: \file{corpo.glb}, \file{coperchio.glb}, \file{filtro.glb}, varie \file{vite\_*.glb}, \file{pavimento.glb}.}
  \caption{La cartella \file{models/} della Pompa Becker: ogni componente smontabile è un file GLB autonomo.}
  \label{fig:cartella-models}
\end{figure}

\subsection{Oggetti compositi con figli interattivi}

Non tutti i modelli devono essere file separati. Quando un oggetto è
un'unità logica con più parti interattive interne, si usa un
\textbf{singolo GLB con figli} (\emph{children}). Nello scenario
Manutenzione Elettromandrino, ad esempio:

\begin{itemize}
  \item \file{pulpito.glb} --- contiene il pannello di controllo con
        tutti i pulsanti (\texttt{Pulsante\_mdi}, \texttt{chiave0},
        \texttt{emergenza\_off}, ecc.) come oggetti figli nella gerarchia
        GLB.
  \item \file{a500.glb} --- contiene l'intera macchina con i
        sotto-assiemi mobili (\texttt{porta}, \texttt{naso}) animabili
        tramite \kw{element = a500.porta}.
  \item \file{remote.glb} --- il telecomando con i pulsanti figli
        (\texttt{pulsante\_r\_play}, \texttt{Pulsante\_r\_unlock}).
\end{itemize}

\begin{notablu}[Aspetti tecnici]
Per ottenere i nomi dei figli di un modello composito, carica la scena
nel browser e usa la console di debug:
\begin{lstlisting}
Scene3D.listChildNames(Scene3D.findModelByName('pulpito'))
\end{lstlisting}
Blender mostra la stessa gerarchia nel pannello \emph{Outliner}.
\end{notablu}

\begin{figure}[H]
  \centering
  % \includegraphics[width=11cm]{immagini/cap02/blender_outliner_pulpito.png}
  \imgph{11cm}{Screenshot del pannello \emph{Outliner} di Blender con la gerarchia di \file{pulpito.glb} espansa: nodo radice \texttt{pulpito} con figli \texttt{Pulsante\_mdi}, \texttt{Pulsante\_tool}, \texttt{chiave0}, \texttt{chiave1}, \texttt{emergenza\_off}, \texttt{pstart0}, \texttt{pstart1}.}
  \caption{L'Outliner di Blender mostra la gerarchia interna del modello \file{pulpito.glb}: ogni figlio diventa un elemento interattivo.}
  \label{fig:blender-outliner}
\end{figure}

\section{Origine, direzione e posizione iniziale}

Ogni modello ha tre proprietà spaziali fondamentali che influenzano il
comportamento nel tutorial.

\subsection{Origine (pivot)}

L'origine definita in Blender è il punto attorno al quale l'oggetto
ruota e rispetto al quale vengono calcolate le traslazioni. Per una
vite, l'origine dovrebbe trovarsi sulla \textbf{testa della vite}
(il punto di presa), non al centro geometrico.

\begin{trucchetto}
In Blender usa \emph{Object $\rightarrow$ Set Origin $\rightarrow$
Origin to 3D Cursor} dopo aver posizionato il cursore 3D nel punto
desiderato. Per le viti, il punto ideale è il centro della testa.
\end{trucchetto}

\begin{figure}[H]
  \centering
  % \includegraphics[width=11cm]{immagini/cap02/origine_vite.png}
  \imgph{11cm}{Confronto a due viti identiche: a sinistra origine al centro geometrico (marcata con una X rossa, con didascalia ``sbagliato''), a destra origine sulla testa della vite (marcata con un punto verde, didascalia ``giusto''). Una freccia curva mostra la rotazione attorno al pivot.}
  \caption{L'origine giusta per una vite: sulla testa, non al centro. Il pivot determina come ruota l'oggetto.}
  \label{fig:origine-vite}
\end{figure}

\subsection{Direzione di movimento}

Ogni modello ha una \textbf{direzione di movimento} associata nel file
\file{scenes/homeconfig.ini}. Questa direzione viene usata da
\val{unscrew}, \val{screw}, \val{extract} e \val{insert} per sapere
lungo quale asse spostare l'oggetto:

\begin{lstlisting}
[models/vite_coperchio_1.glb]
direction = -1, 0, 0    # svitamento verso X negativo

[models/filtro.glb]
direction = -1, 0, 0    # estrazione laterale verso X negativo

[models/vite_culatta_1.glb]
direction = 0, 0, 1     # svitamento verso Z positivo
\end{lstlisting}

\begin{notablu}[Ricorda]
La direzione è un vettore \texttt{(x,~y,~z)} che indica verso dove
l'oggetto si muove quando viene estratto o svitato. Se non viene
specificata, il programma non saprà in che direzione animare l'azione.
\end{notablu}

\begin{figure}[H]
  \centering
  % \includegraphics[width=10cm]{immagini/cap02/sistema_xyz.png}
  \imgph{10cm}{Sistema di assi 3D ortogonali: X rosso (orizzontale), Y verde (verticale), Z blu (profondità). Una vite stilizzata posta nell'origine con freccia che esce lungo $-X$, etichettata ``\texttt{direction = -1, 0, 0}''.}
  \caption{Le tre direzioni dello spazio. Il vettore \texttt{direction} indica dove va l'oggetto quando lo sviti.}
  \label{fig:sistema-xyz}
\end{figure}

\subsection{Posizione iniziale}

La posizione iniziale di ogni modello viene definita in
\file{scenes/homeconfig.ini}. Il sistema crea automaticamente un
riferimento virtuale chiamato \texttt{nomemodello\_original} che
memorizza questa posizione di partenza. Questo riferimento torna utile:

\begin{itemize}
  \item Per riportare un oggetto alla sua posizione iniziale dopo un
        movimento.
  \item Come punto di snap nel Drag \& Drop
        (es.\ \texttt{snap targets filtro\_original}).
  \item Per calcolare traslazioni relative con
        \kw{do : move from filtro\_original by (0,0,0.1)}.
\end{itemize}

%--------------------------------------------------------------------------
\chapter{Organizzare il progetto: cartelle e file}
\label{ch:progetto}
%--------------------------------------------------------------------------

Una volta esportati i modelli 3D, prima di mettere mano al
\file{tutorial.cvtscript} è utile fare il punto sulla struttura del
progetto: dove vanno i file, come si dichiarano gli scenari, quali
controlli effettuare. Sono passaggi rapidi ma fondamentali --- senza una
disposizione corretta dei file, il programma non riesce a trovare nulla.

Questo capitolo è una mini guida ``operativa'' che chiude la fase
preparatoria. I file di configurazione qui introdotti
(\file{homeconfig.ini}, \file{tool.ini}, \file{objects.ini},
\file{interfaceconfig.ini}) hanno una trattazione completa nel
Capitolo~\ref{ch:config}; qui ne diamo solo la panoramica strettamente
necessaria a far girare il tuo primo scenario.

\section{Struttura delle cartelle del progetto}

Le cartelle che contano per scrivere un tutorial sono organizzate in due
livelli. Al \textbf{root del progetto} ci sono le cartelle ``\textbf{globali}''
condivise da tutti gli scenari (modelli 3D, video, immagini per display,
icone della home), mentre dentro \file{scenes/} ci sono i \textbf{file di
configurazione} e una \textbf{sottocartella per ogni scenario}.

\begin{lstlisting}
campus_virtual_training/   # root del progetto
  models/                  # modelli 3D (.glb) — condivisi tra scenari
  media/                   # video e immagini per i message dei tutorial
  screens/                 # frame per le animazioni dei display
  menuimages/              # immagini di anteprima della home page

  scenes/
    homeconfig.ini         # vetrina scenari (home page)
    interfaceconfig.ini    # personalizzazione mouse e camera

    Pompa_Becker/          # un sotto-cartella per ogni scenario
      tutorial.cvtscript   #   sequenza guidata degli step
      tool.ini             #   definizione degli utensili
      objects.ini          #   stati e oggetti interattivi

    Manutenzione_Elettromandrino/
      tutorial.cvtscript
      tool.ini
      objects.ini
\end{lstlisting}

\begin{notablu}[Una volta sola, in alto]
\file{homeconfig.ini} e \file{interfaceconfig.ini} si scrivono \textbf{una
volta sola} dentro \file{scenes/} e valgono per tutti gli scenari. Quando
aggiungi un nuovo scenario, crei una nuova \emph{sottocartella} dentro
\file{scenes/} con i suoi tre file (\file{tutorial.cvtscript},
\file{tool.ini}, \file{objects.ini}) e poi aggiungi una voce
\sez{[Nome scenario]} dentro \file{homeconfig.ini}.
\end{notablu}

\begin{notablu}[Aspetti tecnici]
Il percorso del modello nel tutorial è sempre relativo alla root del
progetto. Scrivendo \texttt{element = vite\_coperchio\_1}, il sistema
cerca automaticamente \file{models/vite\_coperchio\_1.glb} nella cartella
globale al root del progetto. I video citati con \kw{video =} stanno in
\file{media/}, i frame delle animazioni in \file{screens/}.
\end{notablu}

\begin{figure}[H]
  \centering
  % \includegraphics[width=12cm]{immagini/cap02/struttura_cartelle.png}
  \imgph{12cm}{Albero ad icone della struttura del progetto su due livelli: al root le cartelle globali (\file{models/} arancione, \file{media/} viola, \file{screens/} viola, \file{menuimages/} viola) e la cartella \file{scenes/}. Dentro \file{scenes/} i due file \file{.ini} (\file{homeconfig.ini}, \file{interfaceconfig.ini}) e le sottocartelle scenario (\file{Pompa\_Becker/}, \file{Manutenzione\_Elettromandrino/}), ognuna con esattamente tre file: \file{tutorial.cvtscript}, \file{tool.ini}, \file{objects.ini}.}
  \caption{La struttura completa: cartelle globali al root, configurazione e scenari dentro \file{scenes/}.}
  \label{fig:struttura-cartelle}
\end{figure}

\section{Registrare i modelli in \texttt{homeconfig.ini}}

Affinché i modelli vengano caricati all'avvio dello scenario, devono
essere elencati in \file{scenes/homeconfig.ini}. Per ogni modello
servono due righe: il percorso e la direzione di uscita.

\begin{esempiocvt}[Estratto di homeconfig.ini]
\begin{lstlisting}
[Manutenzione pompa del vuoto]
CameraPos    = (-1.5, 0.5, -0.2)
CameraTarget = (0, 0, 0)
tutorial     = scenes/Pompa_Becker/tutorial.cvtscript
tool         = scenes/Pompa_Becker/tool.ini

model     = models/corpo.glb
direction = 0, 0, 1

model     = models/coperchio.glb
direction = -1, 0, 1

model     = models/filtro.glb
direction = -1, 0, 0

model     = models/vite_coperchio_1.glb
direction = -1, 0, 0
# ...altri modelli...
\end{lstlisting}
\end{esempiocvt}

Ogni riga \texttt{model =} dichiara un file GLB da caricare; la riga
\texttt{direction =} immediatamente successiva ne definisce la direzione
di uscita. L'ordine di elencazione non influisce sul risultato finale.

\begin{attenzione}
Se dimentichi di registrare un modello in \file{scenes/homeconfig.ini},
il tutorial potrà riferirsi ad esso con \kw{element =}, ma il modello
non verrà caricato e lo step risulterà vuoto.
\end{attenzione}

\section{Checklist prima di iniziare a scrivere}

Prima di passare alla scrittura del tutorial, verifica che ogni modello
soddisfi questi requisiti:

\begin{enumerate}
  \item[$\square$] Il file è in formato \file{.glb}.
  \item[$\square$] Il nome è in minuscolo, senza spazi, descrittivo.
  \item[$\square$] L'origine è posizionata nel punto logico di
        manipolazione (es.\ testa della vite).
  \item[$\square$] Per i modelli compositi, i figli hanno nomi univoci e
        riconoscibili nell'Outliner di Blender.
  \item[$\square$] I materiali emissivi contengono \texttt{luminoso},
        \texttt{emissive} o \texttt{glow} nel nome.
  \item[$\square$] Il modello è registrato in \file{scenes/homeconfig.ini}
        con la corretta \texttt{direction}.
  \item[$\square$] Il file è copiato nella cartella \file{models/}
        (o nella sottocartella dello scenario).
\end{enumerate}

\begin{trucchetto}
Per verificare che tutti i modelli siano stati caricati correttamente,
avvia lo scenario nel browser e usa il comando di debug
\texttt{Scene3D.listAvailableObjects()} nella console: elencherà tutti
gli oggetti presenti nella scena 3D.
\end{trucchetto}

\begin{figure}[H]
  \centering
  % \includegraphics[width=12cm]{immagini/cap02/checklist_visuale.png}
  \imgph{12cm}{Checklist a 7 voci con icone simboliche per ognuna: file GLB, lettere minuscole, pivot sulla testa di una vite, gerarchia ad albero, etichetta ``glow'' luminosa, registro \file{homeconfig.ini}, cartella \file{models/}. Ogni voce ha una casella $\square$ a sinistra.}
  \caption{Sette punti da spuntare prima di passare alla scrittura del tutorial.}
  \label{fig:checklist-modelli}
\end{figure}

%--------------------------------------------------------------------------
\chapter{Come è organizzato un file \texttt{cvtscript}}
\label{ch:struttura}
%--------------------------------------------------------------------------

Un file \file{tutorial.cvtscript} è composto da \textbf{blocchi}, ciascuno
delimitato da un'intestazione tra parentesi quadre \texttt{[~]}. È la stessa
forma del classico \emph{file INI} usato in molti file di configurazione, con
due peculiarità: i blocchi sono organizzati in una \textbf{gerarchia a tre
livelli} e all'interno di uno step accettano, oltre alle normali coppie
\texttt{chiave = valore}, righe speciali introdotte da \texttt{do :} per
descrivere le animazioni.

\section{La forma di una riga}

Dentro ogni blocco, le righe seguono sempre la stessa forma. Nessuna
eccezione, nessun caso particolare:

\begin{lstlisting}
chiave = valore
\end{lstlisting}

\begin{figure}[H]
  \centering
  % \includegraphics[width=10cm]{immagini/cap03/anatomia_riga.png}
  \imgph{10cm}{Zoom su una singola riga ``\texttt{chiave = valore}'' con tre etichette colorate sopra: ``CHIAVE'' (sopra \texttt{chiave}, blu), ``UGUALE'' (sopra \texttt{=}, arancione), ``VALORE'' (sopra \texttt{valore}, verde). Una freccia che evidenzia gli spazi attorno all'uguale.}
  \caption{Le tre parti di una riga: chiave, uguale, valore. Sempre la stessa struttura, ovunque nel file.}
  \label{fig:anatomia-riga}
\end{figure}

L'unica forma diversa che incontrerai è la riga d'azione \kw{do : verbo \ldots},
trattata nel Capitolo~\ref{ch:do}: il separatore è \texttt{:} invece di
\texttt{=} per evidenziare che si tratta di un \emph{ordine}, non di una
proprietà.

\section{La gerarchia a tre livelli}

Il file è organizzato come la struttura di un libro: \textbf{una cornice
globale}, dentro la quale stanno \textbf{più capitoli}, ciascuno composto da
\textbf{più passi}. I tre livelli si contengono uno dentro l'altro e
corrispondono a tre tipi di blocco:

\begin{longtable}{@{}p{2.4cm}p{3.6cm}p{8cm}@{}}
\toprule
\textbf{Livello} & \textbf{Blocco} & \textbf{Ruolo} \\
\midrule
1 & \sez{[scene]} & La cornice globale: condizioni iniziali della scena 3D, una sola volta all'inizio del file. \\
2 & \sez{[section "Titolo"]} & Un capitolo del tutorial. Un file ne può contenere uno o molti, in sequenza. \\
3 & \sez{[step "Titolo" flag]} & Il singolo passo che compone una sezione. Ogni sezione ne ha tipicamente da 5 a 30. \\
\bottomrule
\end{longtable}

\begin{figure}[H]
  \centering
  % \includegraphics[width=12cm]{immagini/cap03/gerarchia_blocchi.png}
  \imgph{12cm}{Diagramma a tre livelli verticali: in alto un riquadro grande \sez{[scene]} (cornice della stanza, con icona stanza vuota), al centro tre riquadri \sez{[section "..."]} affiancati (capitoli del libro, con icona libro), in basso tanti riquadri piccoli \sez{[step "..."]} (passi singoli, con icona orma). Frecce di contenimento tra i livelli.}
  \caption{I tre livelli della gerarchia: dalla cornice globale al singolo passo.}
  \label{fig:gerarchia-blocchi}
\end{figure}

\subsection{Livello 1 -- \texttt{[scene]}: la cornice globale}

\sez{[scene]} è il blocco di apertura del file. Vi si dichiarano le condizioni
iniziali condivise da tutto il tutorial:

\begin{itemize}
  \item la posizione di partenza della telecamera (\kw{camera = \ldots});
  \item le posizioni iniziali dei modelli 3D che non si trovano già nel
        loro punto naturale (\kw{position = \ldots}, \kw{rotation = \ldots}).
\end{itemize}

Si scrive \textbf{una sola volta}, all'inizio del file:

\begin{lstlisting}
[scene]
camera = position (-1.5, 0.5, -0.2), target (0, 0, 0), zoom 1.2, fade 2.5
\end{lstlisting}

Significa: ``\emph{Telecamera, mettiti in quel punto, guarda verso il centro,
zoomata 1.2, e quando ti sposti impiega 2.5 secondi}''.

\begin{figure}[H]
  \centering
  % \includegraphics[width=11cm]{immagini/cap03/scena_iniziale.png}
  \imgph{11cm}{Render schematico di una scena 3D vuota con la telecamera (a forma di cinepresa) in primo piano e una linea di mira tratteggiata che punta verso il centro della scena. Coordinate \texttt{position (-1.5, 0.5, -0.2)} e \texttt{target (0, 0, 0)} annotate a fianco.}
  \caption{Il blocco \sez{[scene]} fissa la telecamera all'avvio: posizione, target, zoom.}
  \label{fig:scena-iniziale}
\end{figure}

\subsection{Livello 2 -- \texttt{[section "Titolo"]}: il capitolo}

Una \sez{section} raggruppa una sequenza di step affini sotto un titolo
descrittivo. Esempio: il file della pompa Becker contiene quattro sezioni
che corrispondono a quattro tutorial distinti caricabili dalla home:

\begin{enumerate}
  \item Pulizia Filtro
  \item Riassemblaggio
  \item Sostituzione Palette -- Smontaggio
  \item Sostituzione Palette -- Reverse
\end{enumerate}

Ognuna si apre così:

\begin{lstlisting}
[section "Pulizia Filtro"]
\end{lstlisting}

\begin{attenzione}[Le virgolette sono obbligatorie]
Il titolo va sempre fra \textbf{virgolette} \texttt{"\ldots"} perché contiene
spazi. Senza virgolette il parser interpreterebbe ogni parola del titolo
come una bandierina e ti restituirebbe un errore.
\end{attenzione}

\subsection{Livello 3 -- \texttt{[step "Titolo" flag]}: il singolo passo}

Uno \sez{step} è l'unità minima del tutorial: corrisponde a un singolo
passaggio che lo studente deve compiere (svitare una vite, premere un
pulsante, trascinare un filtro). Esempio:

\begin{lstlisting}
[step "Prima vite" highlight]
\end{lstlisting}

L'intestazione contiene il \textbf{titolo} fra virgolette e, opzionalmente,
una o più \textbf{bandierine} che ne modificano il comportamento (in questo
caso \val{highlight} fa pulsare il fumetto della spiegazione). Le bandierine
sono trattate nel Capitolo~\ref{ch:bandierine}; se non te ne servono, ometti
del tutto la parte dopo il titolo:

\begin{lstlisting}
[step "Prima vite"]
\end{lstlisting}

I campi che si possono scrivere dentro lo step (l'\kw{element}, l'eventuale
\kw{tool}, la \kw{description}, le righe \kw{do :}, ecc.) sono il cuore del
linguaggio e occupano i prossimi capitoli, a partire dal
Capitolo~\ref{ch:righe-magiche}.

\section{Altri blocchi accessori}

Oltre ai tre livelli principali, CVTScript prevede alcuni blocchi
\textbf{dichiarativi} per descrivere componenti riutilizzabili:
\sez{[state nome]} per le varianti visive di un oggetto (es. uno schermo
digitale che mostra pagine diverse) e \sez{[object nome]} per i modelli 3D
con parti interattive (pulsanti del pulpito, leve, schermi).

Per igiene di progetto questi blocchi non si scrivono nel
\file{tutorial.cvtscript}: vivono in un file separato chiamato
\file{objects.ini} che il programma carica automaticamente prima del
tutorial. Il Capitolo~\ref{ch:state-object} ne descrive la sintassi; per ora
basti sapere che \emph{esistono} e che è grazie a loro che, dentro lo step,
puoi scrivere semplicemente \kw{element = pulpito.Pulsante\_mdi} e il
programma sa già di che pulsante si tratta.

%--------------------------------------------------------------------------
\chapter{Dentro uno step}
\label{ch:righe-magiche}
%--------------------------------------------------------------------------

\begin{figure}[H]
  \centering
  % \includegraphics[width=13cm]{immagini/cap05/anatomia_step.png}
  \imgph{13cm}{Blocco \texttt{[step "Prima vite" highlight]} esploso, con frecce colorate che escono da ogni riga (\texttt{element}, \texttt{tool}, \texttt{description}, \texttt{do:}, \texttt{camera}, ecc.) e puntano a etichette che spiegano il ruolo di ognuna in linguaggio naturale.}
  \caption{Tutte le righe di uno step: scegli quali ti servono di volta in volta.}
  \label{fig:anatomia-step}
\end{figure}

Ogni step può contenere queste istruzioni (tutte in \textbf{minuscolo} con \texttt{=}):

\begin{longtable}{@{}p{3.3cm}p{4.8cm}p{6.2cm}@{}}
\toprule
\textbf{Riga} & \textbf{Cosa fa} & \textbf{Esempio} \\
\midrule
\kw{element =} & Il pezzo 3D dello step: vite, porta, pulsante, leva, qualsiasi cosa & \texttt{element = vite\_coperchio\_1} \newline \texttt{element = pulpito.Pulsante\_mdi} \\
\kw{tool =} & Quale attrezzo serve allo studente & \texttt{tool = hand} \\
\kw{description =} & Il testo del fumetto & \texttt{description = Rimuovi la prima vite} \\
\kw{do :} & Un'azione (vedi capitolo 5) & \texttt{do : unscrew} \\
\kw{camera =} & Dove guarda la telecamera & \texttt{camera = target filtro, zoom 0.5} \\
\kw{drag =} & Lo studente trascina con il mouse & \texttt{drag = filtro distance 0.3} \\
\kw{after :} & Cosa succede dopo che lo studente clicca l'\kw{element} & \texttt{after : schermo = schermo002} \\
\kw{message =} & Mostra una finestra con un messaggio & \texttt{message = Attenzione, leggi qui!} \\
\kw{title =} & Il titolo della finestra (con \kw{message}) & \texttt{title = Leggi prima di continuare} \\
\kw{video =} & Mostra un video nella finestra & \texttt{video = media/spiegazione.mp4} \\
\kw{image =} & Mostra un'immagine nella finestra & \texttt{image = scenes/X/foto.jpg} \\
\kw{hold =} & Lo studente prende/tiene/lascia un oggetto & \texttt{hold = pick remote at \ldots} \\
\kw{view =} & Cambia la schermata di un display 3D & \texttt{view = main} \\
\kw{position =} & Sposta un pezzo in un punto preciso & \texttt{position = filtro at (0, 0, 0)} \\
\kw{rotation =} & Ruota un pezzo di tot gradi & \texttt{rotation = filtro by (0, 0, 90)} \\
\kw{companion =} & Un pezzo che si muove insieme a quello principale & \texttt{companion = tubo follows master} \\
\kw{animation =} & Mostra una sequenza di immagini animate 2D & \texttt{animation = folder x at (60\%,60\%) frame 20ms} \\
\bottomrule
\end{longtable}

\begin{regola}
Scrivi le righe \textbf{tutte in minuscolo}. Niente \texttt{Element=}, niente
\texttt{TOOL=}. Solo \texttt{element =} e \texttt{tool =}. È il modo per dire al
programma: ``Sto usando la sintassi nuova v3''.
\end{regola}

%--------------------------------------------------------------------------
\chapter{Gli ordini, l'istruzione \texttt{do :}}
\label{ch:do}
%--------------------------------------------------------------------------

\kw{do :} è l'istruzione più importante. Significa \textbf{``fai questa cosa
adesso''}. La riga \kw{do :} può apparire \textbf{più volte} in uno step. Le
azioni vengono fatte \textbf{una dopo l'altra}, nell'ordine in cui le scrivi.

\section{Quando parte un \texttt{do :} ?}

\begin{regola}[La cosa più importante da capire]
Le azioni \kw{do : \ldots} \textbf{non partono da sole all'inizio dello step}.
Partono \textbf{quando lo studente clicca l'\kw{element}} dello step. È
l'interazione dello studente che mette in moto l'animazione.
\end{regola}

In altre parole: ogni step ha un protagonista (l'\kw{element}) e una serie di
cose che gli devono accadere (i \kw{do :}). Lo studente sceglie il momento:
\textbf{cliccando sul protagonista} dà il via alla scena.

\textbf{Esempio.} Vediamo lo step della porta:

\begin{lstlisting}
[step "Apri la porta"]
element = a500.porta
do      : rotate around (0.615, 0, 0) by (0, 90, 0) duration 1.0
\end{lstlisting}

Cosa succede dal punto di vista dello studente:

\begin{enumerate}
  \item Lo step si carica. La porta (\texttt{a500.porta}) si \textbf{illumina
        di giallo} per dire ``cliccami''.
  \item Lo studente \textbf{clicca la porta}.
  \item \emph{Solo a quel punto} parte il \kw{do : rotate \ldots}: la porta
        ruota di 90 gradi in 1 secondo.
  \item Quando l'animazione è finita, lo step avanza al prossimo.
\end{enumerate}

Senza il click dello studente, la porta resta lì ferma. Il programma aspetta.

\subsection{Le tre eccezioni: \texttt{auto}, \texttt{machine}, e l'attrezzo}

Ci sono tre situazioni in cui il \kw{do :} parte \textbf{senza un click
diretto} sull'\kw{element}:

\begin{itemize}
  \item \textbf{Step con bandierina \val{auto}}: lo studente \textbf{ha già
        equipaggiato il tool} dichiarato e il programma simula il click in
        automatico. Si usa per le viti consecutive: la prima la fa lo
        studente, le altre sono in cascata.
  \item \textbf{Step con bandierina \val{machine}}: la macchina si muove da
        sola, senza nemmeno il tool. Esempio: il carro dell'asse Y che si
        sposta automaticamente durante un cambio utensile.
  \item \textbf{Step con \kw{tool = \ldots} + \kw{do :}}: lo studente
        \textbf{equipaggia l'attrezzo dalla cassetta} e poi clicca
        l'\kw{element} con quell'attrezzo. È sempre un click sull'element,
        solo mediato dal tool.
\end{itemize}

\begin{figure}[H]
  \centering
  % \includegraphics[width=13cm]{immagini/cap06/do_quando_parte.png}
  \imgph{13cm}{Quattro vignette in fila. (1) Cursore che scende verso un'\kw{element} illuminata di giallo $\rightarrow$ click $\rightarrow$ \kw{do} che parte. (2) Lo stesso ma il cursore ha forma dell'attrezzo (chiave a brugola). (3) Step con la bandierina \val{auto}: il cursore non c'è, una freccia ``automatico'' avvia il \kw{do}. (4) Step \val{machine}: la macchina che si muove da sola, etichetta ``nessun click necessario''.}
  \caption{Quando parte un \kw{do :}: nei tre casi normali al click dello studente, nei due automatici (\val{auto}/\val{machine}) da solo.}
  \label{fig:do-quando-parte}
\end{figure}

\subsection{Più \texttt{do :} in fila = una sola animazione}

Se metti più righe \kw{do :} nello stesso step, vengono eseguite
\textbf{una dopo l'altra} a partire dal singolo click dello studente. Lo
studente clicca \textbf{una volta}, e parte tutta la sequenza.

\begin{lstlisting}
[step "Estrai la flangia" highlight]
element     = flangia
description = Stacca la flangia dalla macchina.
do          : move by (0, 0, 0.5) duration 0.5
do          : rotate by (90, 0, 0) duration 0.5
do          : place duration 0.2
\end{lstlisting}

\textbf{Cosa vede lo studente.} Clicca la flangia una sola volta. Poi:
prima la flangia si sposta lungo Z (mezzo secondo); poi ruota di 90 gradi
(altro mezzo secondo); poi si appoggia a terra (un quinto di secondo). Tre
animazioni concatenate, un solo click.

\section{Il vocabolario degli ordini}

\begin{figure}[H]
  \centering
  % \includegraphics[width=13cm]{immagini/cap06/9_azioni.png}
  \imgph{13cm}{Tabella visiva 3$\times$3 con le 9 azioni: ognuna ha un'icona simbolica e il nome inglese sotto. \texttt{unscrew} (vite con frecce circolari + uscita), \texttt{screw} (vite che entra), \texttt{extract} (freccia in linea), \texttt{insert} (freccia inversa), \texttt{place} (oggetto che cade), \texttt{move} (cubo con vettore), \texttt{rotate} (cerchio rotatorio), \texttt{pump} (oscillazione $\leftrightarrow$), \texttt{idle} (cerchio reset).}
  \caption{Il vocabolario completo: nove ordini per descrivere ogni movimento.}
  \label{fig:9-azioni}
\end{figure}

Sono solo \textbf{nove parole}:

\begin{longtable}{@{}lp{6cm}p{6cm}@{}}
\toprule
\textbf{Ordine} & \textbf{Cosa fa} & \textbf{Esempio} \\
\midrule
\val{unscrew} & Svita la vite (rotazione + estrazione) & \texttt{do : unscrew distance 0.3} \\
\val{screw} & Avvita la vite (inverso di unscrew) & \texttt{do : screw distance 0.005} \\
\val{extract} & Estrae il pezzo lungo l'asse & \texttt{do : extract distance 0.4} \\
\val{insert} & Inserisce il pezzo & \texttt{do : insert distance 0.4} \\
\val{place} & Appoggia a terra il pezzo & \texttt{do : place duration 0.5} \\
\val{move} & Sposta il pezzo (traslazione) & \texttt{do : move by (0,0,-0.1) duration 0.2} \\
\val{rotate} & Ruota il pezzo & \texttt{do : rotate by (0,0,-90) duration 0.8} \\
\val{pump} & Pompa avanti-indietro (cicli) & \texttt{do : pump along x amplitude 0.08\ldots} \\
\val{idle} & Operazioni speciali & \texttt{do : idle reset} \\
\bottomrule
\end{longtable}

\section{Spiegazioni semplici di ognuno}

\subsection{\texttt{unscrew} -- Svita la vite}

\begin{lstlisting}
do : unscrew                  # Svita di mezzo metro (default)
do : unscrew distance 0.3     # Svita di 30 centimetri
\end{lstlisting}

Pensa a una vite: per svitarla devi \textbf{girarla tante volte} e
\textbf{tirarla fuori}. \val{unscrew} fa entrambe le cose insieme.

\begin{notablu}
La direzione (su, giù, verso destra) è scritta nel file
\file{scenes/homeconfig.ini} per ogni modello. Tu non te ne devi preoccupare.
\end{notablu}

\begin{figure}[H]
  \centering
  % \includegraphics[width=11cm]{immagini/cap06/unscrew_sequenza.png}
  \imgph{11cm}{Tre fotogrammi in sequenza orizzontale: vite avvitata nel pezzo (frame 1), vite a metà strada con freccia di rotazione e freccia di traslazione (frame 2), vite completamente fuori e separata (frame 3). Frecce di flusso tra i frame.}
  \caption{\val{unscrew} combina rotazione e traslazione lungo l'asse di svitatura.}
  \label{fig:unscrew-sequenza}
\end{figure}

\subsection{\texttt{screw} -- Avvita la vite}

L'opposto di \val{unscrew}. La vite \textbf{rientra} nel suo buco.

\begin{lstlisting}
do : screw distance 0.02
\end{lstlisting}

\subsection{\texttt{extract} / \texttt{insert} -- Estrae / Inserisce}

Come svitare ma \textbf{senza la rotazione}. Solo lo spostamento.

\begin{lstlisting}
do : extract distance 0.6     # tira fuori di 60 cm
do : insert  distance 0.6     # rimette dentro
\end{lstlisting}

\begin{figure}[H]
  \centering
  % \includegraphics[width=10cm]{immagini/cap06/extract_insert.png}
  \imgph{10cm}{Pezzo cilindrico (es.\ filtro) che esce in linea retta dalla sua sede, senza rotazione. Sopra l'asse di estrazione una doppia freccia: $\rightarrow$ etichettata \texttt{extract}, $\leftarrow$ etichettata \texttt{insert}. Il pezzo non ruota mai.}
  \caption{\val{extract} è solo traslazione: niente rotazione, solo movimento lungo la direzione del modello.}
  \label{fig:extract-insert}
\end{figure}

\subsection{\texttt{place} -- Appoggia a terra}

Il pezzo \textbf{scende e si appoggia per terra}. Utile dopo che lo hai svitato e tolto.

\begin{lstlisting}
do : place duration 0.5    # impiega mezzo secondo a scendere
\end{lstlisting}

\begin{figure}[H]
  \centering
  % \includegraphics[width=9cm]{immagini/cap06/place.png}
  \imgph{9cm}{Vite sospesa in aria con una freccia verticale verso il basso. Linea tratteggiata orizzontale a Y=0 etichettata ``pavimento''. La vite finisce esattamente sulla linea, appoggiata in piedi.}
  \caption{\val{place} calcola il pavimento e fa scendere il pezzo dolcemente fino a Y=0.}
  \label{fig:place}
\end{figure}

\subsection{\texttt{move} -- Sposta in linea retta}

\begin{lstlisting}
do : move by (1, 0, 0) duration 1.0
\end{lstlisting}

Significa: sposta il pezzo di 1 metro lungo l'asse X, in 1 secondo.

I tre numeri tra parentesi sono \textbf{(X, Y, Z)} -- esattamente come tre direzioni:

\begin{itemize}
  \item \textbf{X} = destra/sinistra (numeri positivi = destra)
  \item \textbf{Y} = su/giù (positivi = su)
  \item \textbf{Z} = avanti/indietro (positivi = avanti)
\end{itemize}

\begin{figure}[H]
  \centering
  % \includegraphics[width=10cm]{immagini/cap06/move_xyz.png}
  \imgph{10cm}{Sistema di riferimento 3D con tre frecce colorate: X rossa (destra/sinistra), Y verde (su/giù), Z blu (avanti/indietro). Un cubo centrato nell'origine con vettore di spostamento \texttt{(1,0,0)} disegnato verso destra.}
  \caption{Le tre coordinate (X,~Y,~Z) come tre direzioni dello spazio.}
  \label{fig:move-xyz}
\end{figure}

\textbf{Versione speciale} che parte da un altro pezzo:

\begin{lstlisting}
do : move from tappino_grasso_dx_original by (-0.2, 0, 0) duration 0.5
\end{lstlisting}

Significa: ``Vai vicino al tappino di destra (alla sua posizione originale),
poi spostati di -0.2 lungo X''. Si usa quando vuoi posizionare un oggetto
\textbf{vicino a un altro}.

\subsection{\texttt{rotate} -- Gira}

\begin{lstlisting}
do : rotate by (0, 0, 90) duration 1.0
\end{lstlisting}

Significa: ruota di 90 gradi attorno all'asse Z, in 1 secondo.

\textbf{Versione speciale} ``ruota intorno a un punto preciso'':

\begin{lstlisting}
do : rotate around (0, 0.1, 0) by (0, 0, -90) duration 0.8
\end{lstlisting}

Significa: ruota di -90 gradi sull'asse Z, ma il \textbf{perno della rotazione}
è il punto (0, 0.1, 0). È come quando una porta gira intorno alla cerniera.

\begin{figure}[H]
  \centering
  % \includegraphics[width=12cm]{immagini/cap06/rotate_due_modi.png}
  \imgph{12cm}{Confronto a due immagini: a sinistra cubo che ruota attorno al suo centro (\texttt{rotate by}, perno = origine pezzo). A destra cubo che ruota attorno a un punto esterno marcato con un cardine (\texttt{rotate around}, come una porta che gira sui cardini).}
  \caption{La differenza tra \texttt{rotate by} (perno = origine pezzo) e \texttt{rotate around} (perno = punto esterno).}
  \label{fig:rotate-due-modi}
\end{figure}

\subsection{\texttt{pump} -- Pompare}

Il pezzo si muove \textbf{avanti e indietro} tante volte (es. una pompetta del grasso).

\begin{lstlisting}
do : pump along x amplitude 0.08 cycles 10 duration 0.1 from tappino_grasso_dx_original
\end{lstlisting}

Si legge: ``Pompa lungo X, ampiezza 8 cm, fai 10 cicli, ogni ciclo dura
0.1 secondi, partendo dal tappino di destra''.

\textbf{Versione asimmetrica} (il ``su'' e il ``giù'' non sono uguali):

\begin{lstlisting}
do : pump from tappino_grasso_dx_original between (-0.08, 0, 0) and (-0.09, 0, 0) cycles 10 duration 0.1
\end{lstlisting}

\begin{figure}[H]
  \centering
  % \includegraphics[width=11cm]{immagini/cap06/pump.png}
  \imgph{11cm}{Pompetta del grasso con stantuffo che si muove avanti-indietro: doppia freccia $\leftrightarrow$ etichettata ``ampiezza''. Sotto, una linea ondulata sinusoidale con N picchi etichettati ``cicli=10''. A fianco una clessidra con scritto ``durata=0.1s per ciclo''.}
  \caption{\val{pump} ripete avanti-indietro N volte: utile per simulare pompette del grasso.}
  \label{fig:pump}
\end{figure}

\subsection{\texttt{idle reset} -- Resetta tutto}

\begin{lstlisting}
do : idle reset
\end{lstlisting}

Mette tutti i pezzi nella loro posizione originale. Si usa solitamente
all'inizio di una \sez{section} per ``ripulire'' la scena.

%--------------------------------------------------------------------------
\chapter{Gli strumenti: i vari attrezzi}
\label{ch:strumenti}
%--------------------------------------------------------------------------

Lo studente, durante il tutorial, sceglie un attrezzo. Tu gli dici quale serve
con la riga \kw{tool = \ldots}. Sono \textbf{poche}, parole inglesi semplici:

\begin{longtable}{@{}llp{8cm}@{}}
\toprule
\textbf{Riga nel file} & \textbf{Attrezzo} & \textbf{Quando usarlo} \\
\midrule
\val{tool = hand}    & Mano                & Per pezzi che si tolgono a mano (tappi, coperchi) \\
\val{tool = hex\_key} & Chiave a brugola    & Per viti a brugola (esagonali) \\
\val{tool = wrench}  & Chiave inglese      & Per dadi e bulloni \\
\val{tool = air}     & Aria compressa      & Per soffiare via la sporcizia \\
\val{tool = spray}   & Spray lubrificante  & Per ingrassare \\
\bottomrule
\end{longtable}

\begin{figure}[H]
  \centering
  % \includegraphics[width=12cm]{immagini/cap07/cinque_attrezzi.png}
  \imgph{12cm}{Cassetta degli attrezzi aperta con cinque scomparti, ognuno con icona e doppia etichetta (italiana sopra, inglese sotto): mano (\texttt{hand}), brugola (\texttt{hex\_key}), chiave inglese (\texttt{wrench}), pistola aria compressa (\texttt{air}), bomboletta spray (\texttt{spray}).}
  \caption{Gli attrezzi disponibili: i nomi inglesi sono quelli da usare nella riga \kw{tool =}.}
  \label{fig:cinque-attrezzi}
\end{figure}

\textbf{Esempio completo:}

\begin{lstlisting}
[step "Rimuovi la vite" highlight]
element     = vite_culatta_1
tool        = hex_key
description = Usa la chiave a brugola da 6mm per svitare la prima vite.
do          : unscrew distance 0.3
do          : place duration 0.2
\end{lstlisting}

\begin{figure}[H]
  \centering
  % \includegraphics[width=12cm]{immagini/cap07/barra_strumenti_app.png}
  \imgph{12cm}{Screenshot della barra degli strumenti reale dell'app: 4 slot in basso allo schermo, ognuno con icona dell'attrezzo e label sotto. Uno slot è evidenziato (selezionato) con bordo colorato. Annotazioni numerate 1--4 sopra ogni slot.}
  \caption{La cassetta degli attrezzi nell'interfaccia: lo studente sceglie da qui prima di ogni azione.}
  \label{fig:barra-strumenti-app}
\end{figure}

\begin{figure}[H]
  \centering
  % \includegraphics[width=12cm]{immagini/cap07/galleria_cursori.png}
  \imgph{12cm}{Galleria di 5 cursori: per ognuno (brugola, chiave inglese, mano, aria, spray) due frame affiancati: ``normale'' a sinistra e ``premuto'' a destra. Etichetta sotto ogni coppia.}
  \caption{I cursori dei cinque attrezzi: cambiano forma quando il mouse è premuto.}
  \label{fig:galleria-cursori}
\end{figure}

%--------------------------------------------------------------------------
\chapter{La camera: dove guarda lo studente}
\label{ch:camera}
%--------------------------------------------------------------------------

La telecamera è gli \textbf{occhi dello studente}. Tu decidi dove guardare con la
riga \kw{camera = \ldots}. Si scrive \textbf{tutto su una riga}, separando con
virgole:

\begin{lstlisting}
camera = position (1, 2, 3), target (0, 0, 0), zoom 1.5, fade 2.0
\end{lstlisting}

\section{Le parti della camera}

\begin{figure}[H]
  \centering
  % \includegraphics[width=12cm]{immagini/cap08/camera_geometria.png}
  \imgph{12cm}{Schema 3D della telecamera: una cinepresa stilizzata in posizione \texttt{position}, vettore tratteggiato verso \texttt{target}, cono giallo del campo visivo (FOV), e cerchio orizzontale che rappresenta il \texttt{pivot}. Misurato anche il \texttt{distance} come segmento tra position e target.}
  \caption{Le proprietà geometriche della camera: ogni numero ha un significato preciso nello spazio.}
  \label{fig:camera-geometria}
\end{figure}

\begin{longtable}{@{}lp{10cm}@{}}
\toprule
\textbf{Pezzettino} & \textbf{Significato} \\
\midrule
\val{position (x, y, z)}      & Dove sta la telecamera nello spazio \\
\val{target (x, y, z)}        & Dove sta guardando \\
\val{target nome\_pezzo}      & Punta a un pezzo (senza coordinate) \\
\val{zoom 1.5}                & Quanto è zoomata (più alto = più vicino) \\
\val{fade 2.0}                & Quanti secondi impiega per spostarsi \\
\val{pivot (x, y, z)}         & Il punto attorno a cui ruota \\
\val{distance 0.5}            & Distanza dal target \\
\val{fov 75}                  & ``Angolo di vista'' della camera \\
\val{rotation (rx, ry, rz)}   & Rotazione della camera \\
\val{free within (-3.14, 3.14)} & Lo studente può girare liberamente \\
\bottomrule
\end{longtable}

\begin{figure}[H]
  \centering
  % \includegraphics[width=12cm]{immagini/cap08/fov_confronto.png}
  \imgph{12cm}{Stessa scena vista con due valori di \texttt{fov}: a sinistra \texttt{fov 45} (campo stretto, effetto teleobiettivo, soggetto grande), a destra \texttt{fov 90} (campo ampio, effetto grandangolo, si vede di più). Cono di vista sovrapposto.}
  \caption{Più alto è il FOV, più ampio è il campo di vista (effetto grandangolo).}
  \label{fig:fov-confronto}
\end{figure}

\begin{figure}[H]
  \centering
  % \includegraphics[width=10cm]{immagini/cap08/free_within.png}
  \imgph{10cm}{Emisfero (cupola) attorno a un oggetto al centro: archi colorati in verde indicano l'intervallo di rotazione consentito da \texttt{free within (-3.14, 3.14)}, archi rossi al di fuori indicano i limiti vietati. Una telecamera si muove lungo l'arco verde.}
  \caption{I limiti \texttt{free within} impediscono allo studente di disorientarsi girando troppo la camera.}
  \label{fig:free-within}
\end{figure}

\begin{trucchetto}
Nelle \sez{section} e negli \sez{step} puoi scrivere \textbf{solo i pezzi che
cambiano}. Il programma ricorda gli altri dal contesto precedente.

Esempio: se uno step è quasi uguale al precedente ma cambi solo il target:

\begin{lstlisting}
[step "Tappino Sinistro" highlight]
camera      = target tappino_grasso_sx
description = Monta il tappino di sinistra.
\end{lstlisting}
\end{trucchetto}

%--------------------------------------------------------------------------
\chapter{Trascinare gli oggetti (drag \& drop)}
\label{ch:drag}
%--------------------------------------------------------------------------

A volte vuoi che sia lo \textbf{studente a trascinare} un pezzo nella sua sede
(non automatico). Si fa con la riga \kw{drag = \ldots}:

\begin{lstlisting}
drag = filtro distance 0.3
\end{lstlisting}

Significa: ``Lo studente può trascinare il filtro. Quando si avvicina entro
30 cm dal punto di destinazione, il pezzo \textbf{si aggancia automaticamente}''.

\begin{figure}[H]
  \centering
  % \includegraphics[width=13cm]{immagini/cap09/snap_sequenza.png}
  \imgph{13cm}{Tre fotogrammi orizzontali: 1) filtro lontano dalla sua sede, cursore mano che lo trascina; 2) filtro entro la sfera tratteggiata di soglia (\texttt{distance 0.3}), il bordo della sfera è visibile; 3) filtro agganciato perfettamente nella sede, indicatore di snap acceso (verde).}
  \caption{La meccanica dello snap: quando l'oggetto entra nella soglia di distanza, si aggancia automaticamente.}
  \label{fig:snap-sequenza}
\end{figure}

\section{Più oggetti contemporaneamente}

\begin{lstlisting}
drag = vite_1, vite_2, vite_3, vite_4 distance 0.2
\end{lstlisting}

Significa: ``Lo studente può trascinare le 4 viti, in qualsiasi ordine''.

\section{Punto di aggancio personalizzato}

\begin{longtable}{@{}lp{9cm}@{}}
\toprule
\textbf{Forma} & \textbf{Cosa fa} \\
\midrule
\val{snap offset (0, 0, 0.01)}      & Si aggancia un po' sollevato (1 cm sopra) \\
\val{snap pivot (x, y, z)}          & Si aggancia con il pivot in quel punto \\
\val{snap targets a, b, c}          & Si aggancia in uno qualsiasi di questi punti \\
\bottomrule
\end{longtable}

\textbf{Esempio con tutto:}

\begin{lstlisting}
drag = filtro distance 1.5 snap offset (0, 0, 0.01)
\end{lstlisting}

\begin{lstlisting}
drag = vite distance 1.5 snap targets foro_1.origin, foro_2.origin
\end{lstlisting}

\begin{figure}[H]
  \centering
  % \includegraphics[width=13cm]{immagini/cap09/snap_varianti.png}
  \imgph{13cm}{Tre riquadri affiancati con esempi di snap: 1) \texttt{snap offset (0,0,0.01)} --- pezzo agganciato 1\,cm sopra la sede (linea tratteggiata di offset visibile); 2) \texttt{snap pivot (x,y,z)} --- pezzo con un perno colorato spostato dal centro; 3) \texttt{snap targets a, b, c} --- vite con 3 fori candidati (bersagli) di cui quello più vicino si illumina.}
  \caption{Tre modi di definire dove un oggetto si aggancia: offset, pivot custom, lista di target.}
  \label{fig:snap-varianti}
\end{figure}

%--------------------------------------------------------------------------
\chapter{Pulsanti, porte, leve: tutto è \texttt{element}}
\label{ch:pulsanti}
%--------------------------------------------------------------------------

\begin{figure}[H]
  \centering
  % \includegraphics[width=12cm]{immagini/cap10/pulpito_annotato.png}
  \imgph{12cm}{Render del pulpito 3D dell'Elettromandrino con callout numerati che evidenziano i singoli pulsanti: \texttt{Pulsante\_mdi}, \texttt{Pulsante\_tool}, \texttt{chiave0}, \texttt{chiave1}, \texttt{emergenza\_off}, \texttt{pstart0}, \texttt{pstart1}. Sopra il modello l'etichetta ``\file{pulpito.glb}''.}
  \caption{Il pulpito di una macchina utensile: ogni pulsante è un figlio del modello \file{pulpito.glb}.}
  \label{fig:pulpito-annotato}
\end{figure}

Una delle cose più importanti del CVTScript v3 è una piccola idea che cambia
tutto: \textbf{tutto quello su cui lo studente clicca si chiama \kw{element}}.
Una vite, un pulsante del pulpito, una porta, una leva, una chiave da girare:
una sola parola, \kw{element}, per dire al programma ``\emph{questo è il pezzo
di questo step}''.

\begin{regola}[La grande regola del v3]
Il programma capisce \textbf{da solo} cosa deve succedere quando lo studente
clicca un \kw{element}, guardando le \emph{altre righe} dello step:
se c'è \kw{do :} parte un'animazione, se c'è \kw{after :} cambia uno stato,
se c'è \kw{tool =} si usa un attrezzo, se c'è la bandierina \val{auto} non
serve neanche cliccare. \textbf{Tu scrivi solo cosa vuoi che succeda}: il
programma collega da sé i pezzi.
\end{regola}

\begin{notablu}[Perché ``element'' e basta]
Nelle prime versioni del v3 c'era una parola \kw{button =} per i pulsanti
3D. Ma una porta non è un pulsante. Una leva non è un pulsante. Una vite
non è un pulsante. Chiamarli tutti ``button'' confondeva chi leggeva il
file. Allora abbiamo unificato: \textbf{tutto è \kw{element}}, e il
contesto chiarisce di cosa si tratta. La parola \kw{button =} continua a
funzionare per non rompere i tutorial vecchi, ma il programma ti avvisa
che è deprecata.
\end{notablu}

\section{Caso A -- Click su un pulsante che cambia uno schermo}

Quando lo studente deve \textbf{premere un pulsante 3D} (es.\ il pulsante MDI
sul pulpito della macchina), scrivi semplicemente l'\kw{element} e il
\kw{after :}:

\begin{lstlisting}
[step "Vai alla schermata MDI" highlight]
description = Premi il pulsante MDI sul pulpito.
element     = pulpito.Pulsante_mdi
after       : schermo = schermo002
\end{lstlisting}

Si legge:

\begin{itemize}
  \item ``Il pezzo dello step è il pulsante chiamato \texttt{Pulsante\_mdi}
        dentro l'oggetto \texttt{pulpito}''.
  \item ``Quando lo studente lo clicca, cambia la variante dello schermo a
        \texttt{schermo002}''.
  \item Niente \kw{do :} significa che non c'è animazione da aspettare: si
        passa subito al prossimo step.
\end{itemize}

\section{Caso B -- Click su una porta che si apre}

Una porta non è un pulsante: è una porta. Ma per il programma è anch'essa
un \kw{element}. La differenza è che ha un'animazione (\kw{do :}) invece di
un \kw{after :}:

\begin{lstlisting}
[step "Apri la porta"]
description = Apri la porta della cabina.
element     = a500.porta
do          : rotate around (0.615, 0, 0) by (0, 90, 0) duration 1.0
\end{lstlisting}

Si legge:

\begin{itemize}
  \item ``Il pezzo dello step è la porta della macchina (\texttt{a500.porta})''.
  \item ``Quando lo studente la clicca, la porta ruota di 90 gradi attorno al
        punto (0.615, 0, 0) in 1 secondo''.
  \item Quando l'animazione è finita, lo step avanza da solo.
\end{itemize}

\section{Caso C -- Click su una vite con un attrezzo}

Quando l'\kw{element} richiede un attrezzo (chiave a brugola, chiave inglese,
spray\ldots) aggiungi \kw{tool = \ldots}. Lo studente prima equipaggia
l'attrezzo, poi clicca l'\kw{element}:

\begin{lstlisting}
[step "Rimuovi vite culatta"]
element     = vite_culatta_1
tool        = hex_key
description = Svita la vite con la chiave a brugola.
do          : unscrew distance 0.3
\end{lstlisting}

\section{\texttt{after} con effetti multipli}

Se cliccare l'\kw{element} deve cambiare \textbf{più cose} contemporaneamente,
separa con \texttt{;}:

\begin{lstlisting}
element = pulpito.pstart0
after   : start = pstart1; schermo = schermo002
\end{lstlisting}

Significa: ``Quando lo studente clicca il pulsante di accensione, cambia sia la
variante del pulsante stesso (acceso) sia quella dello schermo''.

\begin{notablu}[\texttt{after} funziona anche senza click]
\kw{after :} significa ``\textbf{alla fine di questo step}''. Se lo step ha un
\kw{element =} cliccabile, ``alla fine'' = ``dopo il click''. Se invece lo step
è \val{auto} o \val{machine} (senza click), ``alla fine'' = ``quando lo step
parte''.

\textbf{Esempio} -- step \val{machine} che cambia variante senza click:

\begin{lstlisting}
[step "Scambio Utensile" machine]
after : tool = tool0
\end{lstlisting}

Lo studente non deve cliccare niente: la macchina cambia da sola la variante
del tool e va al prossimo step.
\end{notablu}

\section{La tabella che decide tutto}

Ecco la regola completa. Il programma guarda quali righe ci sono nello step e
sceglie \emph{cosa} far succedere al click:

\begin{longtable}{@{}p{2.4cm}p{2.4cm}p{2.4cm}p{2.4cm}p{4.0cm}@{}}
\toprule
\textbf{\kw{tool}?} & \textbf{\kw{do :}?} & \textbf{\kw{after :}?} & \textbf{flag \val{auto}/\val{machine}?} & \textbf{Cosa succede al click} \\
\midrule
sì & sì & --- & no & Lo studente equipaggia il tool, clicca l'\kw{element} con il tool, parte l'animazione. \\
no & sì & --- & no & Lo studente clicca l'\kw{element}, parte l'animazione. \\
qualunque & no & sì & no & Lo studente clicca l'\kw{element}, partono gli effetti di \kw{after :}, lo step avanza. \\
qualunque & no & no & no & Lo studente clicca l'\kw{element}, lo step avanza. \\
qualunque & qualunque & qualunque & sì & Non serve cliccare: parte tutto da solo. L'\kw{element} è solo il bersaglio dell'animazione. \\
\bottomrule
\end{longtable}

\begin{figure}[H]
  \centering
  % \includegraphics[width=13cm]{immagini/cap10/element_unificato.png}
  \imgph{13cm}{Quattro mini-scene affiancate che condividono lo stesso schema di partenza ``il programma guarda lo step''. Scena 1: vite + attrezzo $\to$ ``flusso tool''. Scena 2: porta che ruota $\to$ ``click + animazione''. Scena 3: pulsante che cambia schermo $\to$ ``click + after''. Scena 4: ingranaggio rotante $\to$ ``machine, automatico''. Sopra le scene un'unica scritta: \kw{element = X}.}
  \caption{Una sola keyword, quattro comportamenti: il programma deduce dal contesto.}
  \label{fig:element-unificato}
\end{figure}

\section{Click ripetuti: \texttt{repeat N}}

A volte vuoi che lo studente clicchi lo stesso \kw{element} \textbf{più volte}
(per esempio per simulare ``apri e chiudi la pinza 8 volte e verifica che
funzioni''). Aggiungi \val{repeat N} al valore dell'\kw{element}:

\begin{lstlisting}
[step "Apri e chiudi la pinza" highlight]
element   = remote.Pulsante_r_unlock repeat 8
animation = folder screens/mandrino at (60%, 60%) frame 20ms
\end{lstlisting}

Significa: ``Lo studente deve premere il pulsante \textbf{8 volte}. Ogni volta
appare un'\textbf{immagine diversa} dalla cartella \texttt{screens/mandrino}
(animazione a frame, 20 millisecondi tra una e l'altra). Dopo l'ottavo click,
lo step avanza''.

\begin{figure}[H]
  \centering
  % \includegraphics[width=12cm]{immagini/cap10/repeat_animation.png}
  \imgph{12cm}{Pulsante 3D al centro con contatore ``+1 +1 +1'' che sale (1/8, 2/8\ldots 8/8). Sopra, una sequenza di mini-schermate che cambiano a ogni click (filmstrip). All'ottavo click, il fumetto si chiude e parte la freccia di avanzamento.}
  \caption{\texttt{repeat N} accoppiato a \texttt{animation}: ogni click avanza un frame, all'ennesimo lo step finisce.}
  \label{fig:repeat-animation}
\end{figure}

%--------------------------------------------------------------------------
\chapter{Le bandierine \texttt{auto}, \texttt{machine}, \texttt{highlight}}
\label{ch:bandierine}
%--------------------------------------------------------------------------

\begin{figure}[H]
  \centering
  % \includegraphics[width=12cm]{immagini/cap11/tre_bandierine.png}
  \imgph{12cm}{Tre bandierine stilizzate affiancate, ognuna con icona simbolica e nome: \texttt{auto} (ingranaggio in movimento), \texttt{machine} (silhouette di macchina utensile), \texttt{highlight} (lampada accesa con onde di luce). Sotto ogni bandierina una frase breve ``chi muove cosa, chi attira l'attenzione''.}
  \caption{Le tre bandierine: chi muove cosa, e chi attira l'attenzione dello studente.}
  \label{fig:tre-bandierine}
\end{figure}

Sono \textbf{etichette speciali} che si scrivono \textbf{dopo il nome dello step},
nella parentesi quadra:

\begin{lstlisting}
[step "Nome dello step" auto, highlight]
\end{lstlisting}

\begin{longtable}{@{}lp{11cm}@{}}
\toprule
\textbf{Bandierina} & \textbf{Significato} \\
\midrule
\val{auto} & Lo step si esegue \textbf{da solo}: equipaggia l'attrezzo, fa le azioni,
             e va al prossimo step. Lo studente non deve cliccare. \\
\val{machine} & Come \val{auto} ma \textbf{senza attrezzo}. Si usa per movimenti
                automatici della macchina (es.\ ``il carro Y si sposta''). \\
\val{highlight} & Il fumetto della spiegazione \textbf{lampeggia} per attirare
                  l'attenzione. Si usa per i passi importanti dove vuoi che lo
                  studente legga. \\
\bottomrule
\end{longtable}

Si possono \textbf{combinare} separandole con virgola:

\begin{lstlisting}
[step "Apertura magazzino" highlight, machine]
\end{lstlisting}

Significa: il fumetto lampeggia E lo step si esegue automaticamente.

\begin{figure}[H]
  \centering
  % \includegraphics[width=12cm]{immagini/cap11/highlight_confronto.png}
  \imgph{12cm}{Due screenshot affiancati della stessa scena: a sinistra fumetto normale (testo statico, bordo sottile); a destra fumetto in stato \val{highlight} (bordo luminoso, alone giallo-arancio attorno al riquadro, frecce di pulsazione). Etichetta sotto: ``normale'' / ``highlight''.}
  \caption{L'effetto \val{highlight}: il fumetto pulsa per dirigere l'attenzione dello studente al testo importante.}
  \label{fig:highlight-confronto}
\end{figure}

%--------------------------------------------------------------------------
\chapter{Tenere oggetti in mano (\texttt{hold})}
\label{ch:hold}
%--------------------------------------------------------------------------

A volte lo studente deve \textbf{prendere un oggetto in mano} (per esempio un
telecomando o una pistola di aria compressa) e poi rimetterlo giù. Si gestisce
con la riga \kw{hold = \ldots}, che ha tre forme.

\begin{figure}[H]
  \centering
  % \includegraphics[width=13cm]{immagini/cap12/hold_ciclo.png}
  \imgph{13cm}{Sequenza a tre stati orizzontali: 1) telecomando posato sul tavolo (etichetta \texttt{pick}), 2) telecomando sospeso davanti alla telecamera in primo piano (etichetta \texttt{held}), 3) telecomando di nuovo sul tavolo (etichetta \texttt{release}). Frecce di transizione tra gli stati con i nomi delle azioni.}
  \caption{Il ciclo di vita di un oggetto in mano: pick $\rightarrow$ held $\rightarrow$ release.}
  \label{fig:hold-ciclo}
\end{figure}

\section{\texttt{hold = pick} -- Prendi in mano un oggetto}

\begin{lstlisting}
hold = pick remote at (-0.25, -0.07, 0.25) facing (0, 5, 0)
\end{lstlisting}

Si legge: ``\textbf{Prendi} l'oggetto chiamato \texttt{remote}, mettilo nella
\textbf{posizione} (--0.25, --0.07, 0.25) rispetto alla telecamera, con la
\textbf{rotazione} (0, 5, 0) (in gradi)''.

L'oggetto seguirà la telecamera come se fosse tenuto davanti agli occhi. Le
parti \texttt{at (\ldots)} e \texttt{facing (\ldots)} sono \textbf{opzionali}:
se non le metti il sistema usa una posizione di default.

\begin{figure}[H]
  \centering
  % \includegraphics[width=11cm]{immagini/cap12/hold_at_facing.png}
  \imgph{11cm}{Telecamera in primo piano (cinepresa stilizzata) con il telecomando sospeso davanti agli occhi. Vettore \texttt{at (-0.25, -0.07, 0.25)} disegnato dalla camera al telecomando. Triedro di rotazione \texttt{facing (0, 5, 0)} che mostra l'inclinazione dell'oggetto (gradi).}
  \caption{\texttt{at} posiziona l'oggetto rispetto alla camera; \texttt{facing} ne imposta l'orientamento.}
  \label{fig:hold-at-facing}
\end{figure}

\section{\texttt{hold = held} -- Lo step richiede l'oggetto già in mano}

\begin{lstlisting}
hold = held remote
\end{lstlisting}

Significa: ``Questo step parte solo se l'oggetto \texttt{remote} è
\textbf{già in mano} dal passo precedente''. Si usa quando hai una sequenza di
step in cui l'oggetto viene preso una volta sola e poi usato in più passi.

\section{\texttt{hold = release} -- Rimetti giù}

\begin{lstlisting}
hold = release
\end{lstlisting}

L'oggetto torna nella sua posizione originale nella scena. Lo studente
\textbf{lo molla}.

\begin{esempiocvt}[{Sequenza completa: prendi, usa, lascia}]
\begin{lstlisting}
[step "Prendi il telecomando" highlight]
description = Prendi il telecomando dal tavolo.
tool        = hand
hold        = pick remote at (-0.25, -0.07, 0.25) facing (0, 5, 0)
view        = main

[step "Premi Start sul telecomando"]
description = Premi il pulsante Start.
hold        = held remote
element     = remote.pulsante_r_play
after       : schermo = schermo004

[step "Rimetti giu il telecomando"]
description = Rimetti il telecomando al suo posto.
hold        = release
\end{lstlisting}
\end{esempiocvt}

%--------------------------------------------------------------------------
\chapter{Mostrare un messaggio sullo schermo}
\label{ch:message}
%--------------------------------------------------------------------------

A volte vuoi mostrare allo studente una \textbf{finestra di avviso} con un
testo, magari un'\textbf{immagine} o un \textbf{video} di spiegazione. Si usa
\kw{message =} accompagnato (in modo opzionale) da \kw{title =}, \kw{image =}
e \kw{video =}.

\begin{figure}[H]
  \centering
  % \includegraphics[width=11cm]{immagini/cap13/modal_completo.png}
  \imgph{11cm}{Screenshot di una finestra modale completa: titolo in alto (``Sicurezza importante''), testo del messaggio (``Indossa i guanti\ldots''), video player con controlli al centro, pulsante \texttt{OK} grande in basso. Sfondo dell'app sfumato dietro.}
  \caption{Una finestra modale completa: titolo, messaggio, video opzionale, OK per chiudere.}
  \label{fig:modal-completo}
\end{figure}

\begin{lstlisting}
[step "Attenzione lubrificante" highlight]
message = Indossa i guanti prima di toccare il grasso.
title   = Sicurezza importante
video   = media/lubelett.mp4
\end{lstlisting}

\textbf{Cosa succede}: appare al centro dello schermo una finestra grande con:
\begin{itemize}
  \item Il \textbf{titolo} ``Sicurezza importante''.
  \item Il \textbf{testo} ``Indossa i guanti\ldots''.
  \item Il \textbf{video} che parte automaticamente con i controlli (play/pause).
  \item Un pulsante \textbf{OK} per chiudere e andare avanti.
\end{itemize}

\begin{notablu}[Combinazioni possibili]
\begin{itemize}
  \item \kw{message =} da solo $\rightarrow$ finestra con solo testo
  \item \kw{message =} + \kw{title =} $\rightarrow$ testo con titolo personalizzato
  \item \kw{message =} + \kw{image =} $\rightarrow$ testo + immagine
  \item \kw{message =} + \kw{video =} $\rightarrow$ testo + video (con player)
\end{itemize}

Se non scrivi \kw{title =}, il titolo della finestra sarà uguale al nome dello
step.
\end{notablu}

\begin{figure}[H]
  \centering
  % \includegraphics[width=13cm]{immagini/cap13/modal_varianti.png}
  \imgph{13cm}{Galleria di 4 mock-up di finestra modale, in fila: 1) solo testo (no titolo no media), 2) testo + titolo personalizzato, 3) testo + immagine al centro, 4) testo + video con controlli player. Etichetta sotto ognuna.}
  \caption{Le quattro combinazioni di \kw{message}: scegli ciò che ti serve.}
  \label{fig:modal-varianti}
\end{figure}

\begin{trucchetto}
Puoi mettere \kw{message =} \textbf{insieme} ad altre azioni dello step (es.\ con
\kw{drag =} o \kw{do :}). La finestra appare \textbf{prima}, lo studente la
chiude con OK, poi parte il resto.
\end{trucchetto}

%--------------------------------------------------------------------------
\chapter{Mettere a posto i pezzi prima di iniziare}
\label{ch:position-rotation}
%--------------------------------------------------------------------------

Quando carichi la scena 3D, ogni pezzo si trova nella sua posizione
``naturale'' (quella in cui è stato modellato in Blender). Ma a volte tu
vuoi che la scena \textbf{parta in un certo modo specifico}: per esempio in
un tutorial di rimontaggio, vuoi che le viti siano già \textbf{tolte e
appoggiate per terra}, in modo che lo studente le trascini al loro posto.

Per questo ci sono \kw{position =} e \kw{rotation =}.

\section{\texttt{position} -- Mettere un pezzo in un punto}

\begin{lstlisting}
position = vite_coperchio_1 at (-0.500, -0.349, 0.000)
\end{lstlisting}

Si legge: ``Metti la \texttt{vite\_coperchio\_1} nel punto (--0.5, --0.35, 0)
rispetto al centro della scena''.

\section{\texttt{rotation} -- Ruotare un pezzo}

\begin{lstlisting}
rotation = coperchio by (0, 0, -90)
\end{lstlisting}

Si legge: ``Ruota il \texttt{coperchio} di --90 gradi sull'asse Z''.

\section{Dove si scrivono?}

\begin{itemize}
  \item Nella \sez{[scene]}: si applicano \textbf{una sola volta} all'inizio del
        tutorial.
  \item In una \sez{[section "\ldots"]}: si applicano all'\textbf{inizio del
        capitolo}, prima del primo step.
  \item In uno \sez{[step "\ldots"]}: si applicano \textbf{quando lo step parte}.
\end{itemize}

\begin{esempiocvt}[Posizionare i pezzi smontati per il rimontaggio]
\begin{lstlisting}
[section "Riassemblaggio"]
camera   = position (-1.5, 0.5, -0.2), target (0, 0, 0), zoom 1.2

# Pezzi che partono smontati e per terra
position = coperchio at (-1.148, -0.261, 0.500)
rotation = coperchio by (0, 0, -90)
position = filtro at (-0.300, -0.142, 0.500)
position = vite_coperchio_1 at (-0.500, -0.349, 0)
position = vite_coperchio_2 at (-0.500, -0.349, 0)

[step "Tappino destro" highlight]
description = Monta il tappino di destra.
tool        = hand
drag        = tappino_grasso_dx distance 0.2
\end{lstlisting}
\end{esempiocvt}

\begin{figure}[H]
  \centering
  % \includegraphics[width=13cm]{immagini/cap14/prima_dopo_riassemblaggio.png}
  \imgph{13cm}{Confronto a due render della stessa scena: a sinistra ``Default'' --- pompa intera già montata; a destra ``Dopo \texttt{position}/\texttt{rotation}'' --- coperchio per terra ruotato di 90$^\circ$, filtro a fianco, viti sparpagliate sul pavimento. Etichetta sopra: ``Stato iniziale per il rimontaggio''.}
  \caption{La stessa scena prima e dopo i comandi \texttt{position}/\texttt{rotation} della sezione di rimontaggio.}
  \label{fig:prima-dopo-riassemblaggio}
\end{figure}

%--------------------------------------------------------------------------
\chapter{Pezzi che si muovono insieme (\texttt{companion})}
\label{ch:companion}
%--------------------------------------------------------------------------

A volte un pezzo si muove e \textbf{ne tira un altro}. Per esempio: quando tiri
fuori la flangia della pompa, il \textbf{tubo del grasso} attaccato dietro
viene tirato anche lui.

In questi casi il pezzo principale (che lo studente sposta) si chiama
\textbf{master}. I pezzi che lo seguono si chiamano \textbf{compagni}
(\val{companion}).

\section{\texttt{companion = X follows master} -- Mi attacco al master}

\begin{lstlisting}
[step "Allontana la flangia"]
element   = flangia
tool      = hand
companion = tubograsso follows master
do        : move by (0, 0, 0.2) duration 0.5
\end{lstlisting}

Significa: ``Quando la \textbf{flangia} (master) si sposta di 20 cm avanti, il
\textbf{tubograsso} la \textbf{segue} attaccato''.

\begin{figure}[H]
  \centering
  % \includegraphics[width=11cm]{immagini/cap15/follows_master.png}
  \imgph{11cm}{Flangia (master, blu) con catena/corda visuale che la lega al tubograsso (companion, arancione). Una freccia di movimento parte dalla flangia, e il tubograsso si muove rigidamente nella stessa direzione e quantità. Etichetta ``master'' sulla flangia, ``companion'' sul tubo.}
  \caption{\val{follows master}: il companion è agganciato e segue rigidamente il movimento del master.}
  \label{fig:follows-master}
\end{figure}

\section{\texttt{companion = X moves by (\ldots) duration t} -- Movimento separato}

\begin{lstlisting}
companion = flangia moves by (0, 0, 0.005) duration 0.5
companion = tubograsso moves by (0, 0, 0.005) duration 0.5
\end{lstlisting}

In questo caso il companion fa un \textbf{movimento suo}, in parallelo al
master. Si usa per esempio quando si avvitano viti estrattori e \textbf{tutta
la flangia avanza un pochino} a ogni giro di vite.

\begin{figure}[H]
  \centering
  % \includegraphics[width=11cm]{immagini/cap15/moves_by.png}
  \imgph{11cm}{Master e companion disegnati come due oggetti distinti, ognuno con la propria freccia di movimento parallela ma indipendente. Una clessidra centrale unisce le due frecce per indicare ``stessa durata''. Nessuna catena tra i due.}
  \caption{\val{moves by}: il companion fa il proprio movimento in parallelo, non agganciato al master.}
  \label{fig:moves-by}
\end{figure}

\begin{notablu}
Si possono mettere \textbf{più righe} \kw{companion = \ldots} nello stesso step:
ogni pezzo segue/si muove come scritto nella sua riga.
\end{notablu}

%--------------------------------------------------------------------------
\chapter{Schermate di un display 3D (\texttt{view})}
\label{ch:view}
%--------------------------------------------------------------------------

\begin{figure}[H]
  \centering
  % \includegraphics[width=13cm]{immagini/cap16/display_viste.png}
  \imgph{13cm}{Quattro pannelli affiancati di uno stesso display 3D: ``home'' (icone delle app), ``mdi'' (campo di input), ``tool'' (parametri utensile), ``running'' (barra avanzamento). Ogni vista ha una etichetta sopra con il nome \texttt{view}.}
  \caption{Le viste di un display digitale: ogni \texttt{view} è un GLB diverso che si scambia in tempo reale.}
  \label{fig:display-viste}
\end{figure}

Alcuni oggetti della scena hanno uno \textbf{schermo digitale} (per esempio il
display di una macchina utensile, o lo schermo di un telecomando). Questo
schermo può mostrare \textbf{varie pagine} -- come il telefono che cambia
applicazione.

Quando in uno step vuoi \textbf{cambiare la pagina} mostrata da un display,
usi:

\begin{lstlisting}
view = main
\end{lstlisting}

Significa: ``Mostra la pagina chiamata \texttt{main} sullo schermo''. Le pagine
disponibili sono definite nel file \file{objects.ini} (vedi capitolo dedicato).

\begin{esempiocvt}[Prendi il telecomando con la sua schermata principale]
\begin{lstlisting}
[step "Prendi il telecomando"]
description = Prendi il telecomando dal tavolo.
tool        = hand
hold        = pick remote at (-0.25, -0.07, 0.25) facing (0, 5, 0)
view        = main
\end{lstlisting}
\end{esempiocvt}

\begin{figure}[H]
  \centering
  % \includegraphics[width=11cm]{immagini/cap16/state_machine_viste.png}
  \imgph{11cm}{Diagramma a stati con tre cerchi nominati ``main'', ``pumping'', ``idle'' collegati da frecce con etichette degli step (es.\ ``view = pumping''). Mostra le transizioni possibili tra le viste di un display.}
  \caption{Le transizioni tra viste: dichiarate in \file{objects.ini}, attivate dal tutorial con \texttt{view =}.}
  \label{fig:state-machine-viste}
\end{figure}

%--------------------------------------------------------------------------
\chapter{Animazioni 2D a frame (\texttt{animation})}
\label{ch:animation}
%--------------------------------------------------------------------------

Per simulare \textbf{contatori, schermate che cambiano, sequenze ripetitive}
puoi mostrare una serie di immagini sopra la scena 3D, una dopo l'altra. Si fa
con:

\begin{figure}[H]
  \centering
  % \includegraphics[width=13cm]{immagini/cap17/filmstrip.png}
  \imgph{13cm}{Filmstrip orizzontale di 8 frame in fila (es.\ contatore numerico che sale 1, 2, 3\ldots oppure pinza che si apre progressivamente). Sotto: timeline con tacche \texttt{20ms} tra ogni frame, e freccia di scorrimento ``$\rightarrow$''.}
  \caption{Un'animazione a frame: una sequenza di immagini mostrate ogni N millisecondi sopra la scena 3D.}
  \label{fig:filmstrip}
\end{figure}

\begin{lstlisting}
animation = folder screens/mandrino at (60%, 60%) frame 20ms
\end{lstlisting}

Significa:
\begin{itemize}
  \item \val{folder screens/mandrino} -- la cartella dove ci sono le immagini
        (le carica tutte in ordine di nome)
  \item \val{at (60\%, 60\%)} -- in che punto dello schermo (percentuali della
        finestra: 0\%=sinistra/alto, 100\%=destra/basso)
  \item \val{frame 20ms} -- quanti millisecondi tra un'immagine e la prossima
\end{itemize}

Tutte e tre le parti sono opzionali (puoi metterne anche solo due).

\begin{figure}[H]
  \centering
  % \includegraphics[width=12cm]{immagini/cap17/animazione_posizione.png}
  \imgph{12cm}{Screenshot della finestra dell'app con un riquadro semi-trasparente sovrapposto alla scena 3D, posizionato in alto-destra (al 60\%, 60\%). Linee guida tratteggiate mostrano le percentuali ``60\%'' dall'angolo in alto a sinistra (X) e dall'alto (Y).}
  \caption{Posizionamento dell'animazione: percentuali sulla finestra (0\%~=~sinistra/alto, 100\%~=~destra/basso), non pixel.}
  \label{fig:animazione-posizione}
\end{figure}

\begin{notablu}[Quasi sempre va con \texttt{element \ldots repeat N}]
L'animazione è quasi sempre \textbf{guidata da clic ripetuti}: ogni clic fa
avanzare di un frame. Nello stesso step si scrive l'\kw{element} con
\val{repeat N}:

\begin{lstlisting}
[step "Apri e chiudi" highlight]
description = Apri e chiudi piu volte premendo il pulsante (4 cicli)
tool        = hand
element     = remote.Pulsante_r_unlock repeat 8
animation   = folder screens/mandrino at (60%, 60%) frame 20ms
\end{lstlisting}

Lo studente preme 8 volte il pulsante, ad ogni pressione l'immagine cambia,
all'ottava lo step finisce.
\end{notablu}

%--------------------------------------------------------------------------
\chapter{I blocchi \texttt{[state]} e \texttt{[object]}}
\label{ch:state-object}
%--------------------------------------------------------------------------

Oltre ai tre blocchi gerarchici (\sez{[scene]}, \sez{[section]},
\sez{[step]}) visti nel Capitolo~\ref{ch:struttura}, esistono due blocchi
\textbf{dichiarativi} che servono a descrivere oggetti con \textbf{più stati
visivi} (per esempio una chiave che può essere ``girata'' o ``non girata'')
oppure modelli 3D \textbf{con parti interattive} (per esempio un pulpito di
comando con i suoi pulsanti).

Per igiene di progetto questi blocchi \textbf{non si scrivono nel}
\file{tutorial.cvtscript}, ma in un file separato chiamato
\file{objects.ini} caricato automaticamente prima del tutorial.

\section{\texttt{[state nome]} -- Le varianti di un oggetto}

Pensa a una \textbf{lampadina}: può essere \textbf{accesa} o \textbf{spenta}.
Sono due \textbf{varianti} dello stesso oggetto. Si scrive così:

\begin{lstlisting}
[state lampadina]
Variants = lampadina_off, lampadina_on
Default  = lampadina_off
\end{lstlisting}

Significa:
\begin{itemize}
  \item Lo state si chiama \texttt{lampadina}
  \item Ha due varianti: \texttt{lampadina\_off} e \texttt{lampadina\_on}
        (sono i nomi dei mesh dentro al modello 3D)
  \item All'inizio è in \texttt{lampadina\_off}
\end{itemize}

In uno step puoi cambiare variante con \kw{after :}:

\begin{lstlisting}
[step "Accendi" machine]
after : lampadina = lampadina_on
\end{lstlisting}

\begin{figure}[H]
  \centering
  % \includegraphics[width=11cm]{immagini/cap18/state_lampadina.png}
  \imgph{11cm}{Due icone di lampadina affiancate: a sinistra spenta (etichetta \texttt{lampadina\_off}), a destra accesa con alone giallo (etichetta \texttt{lampadina\_on}). Una doppia freccia tra le due con scritta ``setVariant''. Sopra il riquadro: ``state lampadina''.}
  \caption{Uno \texttt{state} è un gruppo di varianti mutuamente esclusive: la lampadina è O accesa O spenta, mai entrambe.}
  \label{fig:state-lampadina}
\end{figure}

\section{\texttt{[object nome]} -- Un oggetto con pulsanti dentro}

Quando un modello 3D ha dei \textbf{pulsanti} che lo studente può cliccare
(per esempio il pulpito di una macchina), lo dichiari così:

\begin{lstlisting}
[object pulpito]
Model = models/pulpito.glb

InteractiveChild = Pulsante_mdi, button, onClick:setVariant:schermo=schermo002
InteractiveChild = chiave0, button, onClick:cycleVariant:chiave
\end{lstlisting}

Si legge:
\begin{itemize}
  \item L'oggetto si chiama \texttt{pulpito} e usa il modello
        \texttt{models/pulpito.glb}
  \item Dentro c'è un pulsante \texttt{Pulsante\_mdi}: quando lo clicchi, cambia
        la variante dello state \texttt{schermo} a \texttt{schermo002}.
  \item C'è anche \texttt{chiave0}: quando lo clicchi, cicla
        sulla \textbf{prossima variante} dello state \texttt{chiave} (off
        $\rightarrow$ on $\rightarrow$ off\ldots).
\end{itemize}

\begin{notablu}[A che servono?]
Il file \file{objects.ini} \textbf{descrive una sola volta} che il pulpito ha 5
pulsanti, che lo schermo può mostrare 9 pagine, eccetera. Poi nel tutorial usi
solo \kw{element =} per attivarli e \kw{after :} per cambiare le pagine.

È la stessa cosa che fai con le \textbf{persone}: prima dici ``Mario è alto e ha
i capelli neri'' (= dichiari l'oggetto), poi puoi parlare di Mario senza
ridescriverlo ogni volta.
\end{notablu}

\begin{figure}[H]
  \centering
  % \includegraphics[width=12cm]{immagini/cap18/object_pulpito.png}
  \imgph{12cm}{Render del pulpito 3D con etichette colorate sui figli interattivi: \texttt{Pulsante\_mdi $\rightarrow$ setVariant:schermo=schermo002}, \texttt{chiave0 $\rightarrow$ cycleVariant:chiave}, \texttt{Pulsante\_tool $\rightarrow$ setVariant:schermo=schermo003}. Sopra: titolo ``object pulpito''.}
  \caption{Un \texttt{object} è un modello 3D con figli cliccabili: ogni \texttt{InteractiveChild} ha la sua azione \texttt{onClick}.}
  \label{fig:object-pulpito}
\end{figure}

%--------------------------------------------------------------------------
\chapter{I file di configurazione attorno al tutorial}
\label{ch:config}
%--------------------------------------------------------------------------

\begin{figure}[H]
  \centering
  % \includegraphics[width=13cm]{immagini/cap19/mappa_orbitale.png}
  \imgph{13cm}{Mappa orbitale: al centro un planeta blu con etichetta ``\file{tutorial.cvtscript}''. In orbita 4 satelliti collegati da linee con frecce: \file{homeconfig.ini} (vetrina scenari), \file{interfaceconfig.ini} (mouse e camera), \file{tool.ini} (attrezzi), \file{objects.ini} (stati e pulsanti). Ogni satellite ha sotto una mini-etichetta che spiega il ruolo.}
  \caption{I cinque file che compongono uno scenario, e come si parlano tra loro.}
  \label{fig:mappa-orbitale}
\end{figure}

Il file \file{tutorial.cvtscript} \textbf{non è solo}. Vive insieme a
\textbf{quattro fratellini} che si trovano in altre cartelle. Ognuno ha un
compito preciso:

\begin{longtable}{@{}p{4.2cm}p{8.5cm}@{}}
\toprule
\textbf{File} & \textbf{A cosa serve} \\
\midrule
\file{scenes/homeconfig.ini} & Lista degli scenari disponibili (la pagina iniziale che vedi quando apri il programma) \\
\file{scenes/interfaceconfig.ini} & Personalizzazione dell'interfaccia (zoom, rotazione camera, sensibilità mouse) \\
\file{scenes/<NomeScenario>/tool.ini} & Quali attrezzi vede lo studente per questo scenario \\
\file{scenes/<NomeScenario>/objects.ini} & Stati e oggetti interattivi (pulsanti, schermi) della scena \\
\bottomrule
\end{longtable}

\section{\texttt{homeconfig.ini} -- La vetrina degli scenari}

È la \textbf{lista dei tutorial} che vede lo studente quando apre il programma.
Ogni scenario è un \textbf{blocco}, e dentro dichiari di che tutorial si tratta,
quali modelli 3D caricare, dove mettere la telecamera all'inizio.

\begin{esempiocvt}[Esempio: una pagina con due scenari]
\begin{lstlisting}
title    = Campus Virtual Training Morbidelli X400
subtitle = Seleziona uno scenario per cominciare

[Manutenzione pompa del vuoto]
CameraPos        = (-1.5, 0.5, -0.2)
CameraTarget     = (0, 0, 0)
AmbientLight     = 0x404040, 1.5
DirectionalLight = 0xffffff, 2.5, (-8, 15, 10)
description      = Smontaggio e manutenzione della pompa Becker.
image            = menuimages/pompavuoto.png
tutorial         = scenes/Pompa_Becker/tutorial.cvtscript
tool             = scenes/Pompa_Becker/tool.ini
model            = models/culatta.glb
direction        = 0, 0, 1
model            = models/corpo.glb
direction        = 0, 0, 1

[Manutenzione Elettromandrino]
CameraPos        = (-1.11, 1.44, 3.97)
description      = Pulizia, lubrificazione e blocco/sblocco pinza.
image            = menuimages/1.png
tutorial         = scenes/Manutenzione_Elettromandrino/tutorial.cvtscript
tool             = scenes/Manutenzione_Elettromandrino/tool.ini
model            = models/pulpito.glb
direction        = 0, 0, 1
\end{lstlisting}
\end{esempiocvt}

\textbf{Le righe importanti}:

\begin{longtable}{@{}lp{9cm}@{}}
\toprule
\textbf{Riga} & \textbf{Cosa fa} \\
\midrule
\val{title} / \val{subtitle} & Titolo e sottotitolo della home page \\
\val{[Nome scenario]} & Apre il blocco dello scenario (= card visibile in home) \\
\val{description} & Testo descrittivo dello scenario \\
\val{image} & L'immagine di anteprima della card \\
\val{tutorial} & Il file \file{tutorial.cvtscript} da caricare \\
\val{tool} & Il file \file{tool.ini} da caricare per gli attrezzi \\
\val{CameraPos / CameraTarget} & Dove sta la telecamera all'apertura della scena \\
\val{AmbientLight} & Luce ambiente (colore esadecimale, intensità) \\
\val{DirectionalLight} & Luce principale (colore, intensità, direzione) \\
\val{model} & Un modello 3D da caricare nella scena (può apparire più volte) \\
\val{direction} & Per ogni \texttt{model}, la direzione di svita/estrazione \\
\bottomrule
\end{longtable}

\begin{notablu}[\texttt{direction} a che serve?]
Quando scrivi \texttt{do : unscrew} su una vite, il programma deve sapere
\textbf{in che direzione} esce. La risposta sta in \kw{direction} di quel
modello in \file{homeconfig.ini}: \texttt{(0, 0, 1)} = esce verso Z+,
\texttt{(--1, 0, 0)} = esce verso X--, eccetera. Tu ti preoccupi solo di scrivere
\texttt{direction} una volta per modello.
\end{notablu}

\begin{figure}[H]
  \centering
  % \includegraphics[width=12cm]{immagini/cap19/home_page_app.png}
  \imgph{12cm}{Screenshot della home page reale dell'applicazione: titolo ``Campus Virtual Training Morbidelli X400'' in alto, sottotitolo, due card affiancate (``Manutenzione pompa del vuoto'' con immagine pompa, ``Manutenzione Elettromandrino'' con immagine mandrino), pulsante ``Avvia'' su ognuna.}
  \caption{La home page generata da \file{homeconfig.ini}: ogni \texttt{[Nome scenario]} produce una card.}
  \label{fig:home-page-app}
\end{figure}

\section{\texttt{tool.ini} -- Gli attrezzi della cassetta}

Per ogni scenario decidi \textbf{quali 4 attrezzi} vede lo studente nella barra
in basso. È come scegliere gli attrezzi da mettere nello zainetto prima di
partire.

\begin{esempiocvt}[Esempio: scenario Pompa Becker]
\begin{lstlisting}
[Tools]
Slot1 = brugola
Slot2 = chiave_inglese
Slot3 = mano
Slot4 = aria

[Tool:brugola]
Label              = Brugola
Icon               = utilimages/brugola.png
Cursor             = cursors/brugola_normal.svg
CursorPressed      = cursors/brugola_premuto_frame1.svg
CursorPressedFrame2= cursors/brugola_premuto_frame2.svg
CursorHotspotX     = 4
CursorHotspotY     = 9
Type               = tool
TutorialNames      = ChiaveBrugola, hex_key

[Tool:mano]
Label         = Mano
Icon          = utilimages/mano.png
Cursor        = cursors/mano_normal.svg
CursorPressed = cursors/mano_premuto_frame1.svg
CursorHotspotX= 16
CursorHotspotY= 16
Type          = tool
TutorialNames = Mani, hand
\end{lstlisting}
\end{esempiocvt}

\textbf{Spiegato facile facile}:

\begin{longtable}{@{}lp{10cm}@{}}
\toprule
\textbf{Riga} & \textbf{Cosa fa} \\
\midrule
\sez{[Tools]} & Il blocco con i 4 slot della barra attrezzi \\
\val{Slot1\ldots Slot4} & Quale attrezzo va in ogni posizione (da sinistra a destra) \\
\sez{[Tool:nome]} & Apre il blocco con i dettagli di quell'attrezzo \\
\val{Label} & Il nome che appare in italiano sotto l'icona \\
\val{Icon} & L'immagine che si vede nella barra \\
\val{Cursor} & Come diventa il cursore del mouse quando l'attrezzo è scelto \\
\val{CursorPressed} / \val{CursorPressedFrame2} & I due frame dell'animazione del cursore quando si clicca \\
\val{CursorHotspotX/Y} & Il puntino preciso del cursore (la ``punta'' dell'attrezzo) \\
\val{Type} & \texttt{tool} = è un attrezzo (al momento è l'unico valore) \\
\val{TutorialNames} & I nomi che il tutorial può usare nella riga \kw{tool =} per attivarlo \\
\bottomrule
\end{longtable}

\begin{trucchetto}
\val{TutorialNames} accetta \textbf{più nomi} separati da virgola. Così puoi
scrivere indifferentemente \texttt{tool = hand} o \texttt{tool = Mani} e il
programma sceglie comunque la mano.
\end{trucchetto}

\begin{figure}[H]
  \centering
  % \includegraphics[width=12cm]{immagini/cap19/barra_4_slot.png}
  \imgph{12cm}{Barra strumenti reale dell'app con 4 slot in basso, ognuno etichettato sopra con ``Slot1'', ``Slot2'', ``Slot3'', ``Slot4'' (numerati da sinistra a destra). Sotto ogni slot l'icona dell'attrezzo: brugola, chiave inglese, mano, aria.}
  \caption{I 4 slot della barra strumenti: l'ordine in \file{tool.ini} è l'ordine sullo schermo, da sinistra a destra.}
  \label{fig:barra-4-slot}
\end{figure}

\section{\texttt{objects.ini} -- I pulsanti e gli schermi della scena}

Qui dichiari i \textbf{blocchi dichiarativi} (\sez{[state]} e \sez{[object]}) che
descrivono \textbf{cosa è cliccabile} nella scena 3D e \textbf{come si
comporta} quando viene cliccato.

\begin{esempiocvt}[Esempio: pulpito di una macchina utensile]
\begin{lstlisting}
# La chiave puo essere girata o non girata
[state chiave]
Variants = chiave0, chiave1
Default  = chiave0

# Lo schermo del pulpito puo mostrare diverse pagine
[state schermo]
Variants = schermo001, schermo002, schermo003, schermo004
Default  = schermo001

# Il pulpito ha pulsanti e una chiave dentro
[object pulpito]
Model = models/pulpito.glb

InteractiveChild = chiave0, button, onClick:cycleVariant:chiave
InteractiveChild = chiave1, button, onClick:cycleVariant:chiave
InteractiveChild = Pulsante_mdi, button, onClick:setVariant:schermo=schermo002
InteractiveChild = Pulsante_tool, button, onClick:setVariant:schermo=schermo003
\end{lstlisting}
\end{esempiocvt}

\textbf{Cosa significa}:

\begin{itemize}
  \item Lo \textbf{state \texttt{chiave}} dice ``c'è un mesh che può essere
        \texttt{chiave0} oppure \texttt{chiave1}; all'inizio è
        \texttt{chiave0}''.
  \item Lo \textbf{state \texttt{schermo}} elenca le 4 pagine possibili dello
        schermo del pulpito.
  \item L'\textbf{oggetto \texttt{pulpito}} carica il modello 3D e dichiara che
        ha dei figli interattivi (i pulsanti).
  \item Per ogni pulsante \kw{InteractiveChild} dice: ``Quando lo clicchi, fai
        questa azione''.
\end{itemize}

\textbf{Le azioni possibili al click}:

\begin{longtable}{@{}lp{10cm}@{}}
\toprule
\textbf{Forma} & \textbf{Cosa fa} \\
\midrule
\val{onClick:setVariant:STATE=VARIANTE} & Imposta lo state alla variante data \\
\val{onClick:cycleVariant:STATE} & Va alla variante successiva (cicla) \\
\bottomrule
\end{longtable}

\begin{notablu}[Perché è separato dal tutorial?]
Il pulpito è \textbf{sempre lo stesso} -- ha gli stessi pulsanti -- in tutti i
tutorial. Se descrivessi i pulsanti dentro \file{tutorial.cvtscript}, dovresti
ripetere la stessa cosa in ogni file. Mettendo le definizioni in
\file{objects.ini} le scrivi \textbf{una volta sola}.
\end{notablu}

\section{\texttt{interfaceconfig.ini} -- Personalizzare i comandi}

Questo file decide come si comporta il \textbf{mouse} (e il pinch-to-zoom su
touch). Non serve modificarlo per scrivere un tutorial, ma è utile sapere che
esiste.

\begin{esempiocvt}[Configurazione tipica]
\begin{lstlisting}
[ResetCameraButton]
Position = bottom-right

[CameraControls]
InvertVertical    = true
InvertHorizontal  = false
PinchSensitivity  = 0.001
InvertPinchZoom   = false
ScrollSensitivity = 0.01
InvertScrollZoom  = false
ZoomMin           = 0.15
ZoomMax           = 1.3

[SpeechBubble]
DramaticAnimationDuration = 2000
\end{lstlisting}
\end{esempiocvt}

\textbf{Le impostazioni più importanti}:

\begin{longtable}{@{}lp{10cm}@{}}
\toprule
\textbf{Riga} & \textbf{Cosa fa} \\
\midrule
\val{Position} (in \sez{[ResetCameraButton]}) & Dove sta il pulsantino di reset
camera. Valori: \texttt{top-right}, \texttt{top-left}, \texttt{top-center},
\texttt{bottom-right}, \texttt{bottom-left}, \texttt{bottom-center} \\
\val{InvertVertical} & Inverte la rotazione su/giù della camera (\texttt{true}
= naturale, \texttt{false} = invertito) \\
\val{InvertHorizontal} & Inverte la rotazione sinistra/destra \\
\val{PinchSensitivity} & Quanto è sensibile il pinch-to-zoom su touch
(0.001--0.010, default 0.003) \\
\val{ScrollSensitivity} & Sensibilità della rotellina del mouse
(0.010--0.060) \\
\val{ZoomMin / ZoomMax} & Quanto vicino e quanto lontano si può zoomare \\
\val{DramaticAnimationDuration} & Quanti millisecondi dura l'animazione di
``ingrandimento drammatico'' del fumetto quando ha la bandierina \val{highlight}
\\
\bottomrule
\end{longtable}

\begin{regola}
\textbf{Quasi mai} devi toccare \file{interfaceconfig.ini}. È pensato per
piccoli ritocchi all'esperienza, non per il contenuto del tutorial.
\end{regola}

%--------------------------------------------------------------------------
\chapter{Un esempio completo, riga per riga}
\label{ch:esempio}
%--------------------------------------------------------------------------

Vediamo uno \textbf{step completo} e capiamolo riga per riga.

\begin{esempiocvt}[Step ``Rimuovi la prima vite della culatta'']
\begin{lstlisting}
[step "Rimuovi la prima vite della culatta" highlight]
element     = vite_culatta_1
tool        = hex_key
description = Usa la chiave a brugola da 6mm per svitare la prima vite della culatta.
do          : unscrew distance 0.3
do          : place duration 0.2
\end{lstlisting}
\end{esempiocvt}

\textbf{Cosa succede quando lo studente arriva qui:}

\begin{enumerate}
  \item Apparirà un \textbf{fumetto} con scritto ``Usa la chiave a brugola
        da 6mm\ldots'' (e il fumetto lampeggia, perché c'è la bandierina
        \val{highlight}).
  \item Lo studente deve \textbf{scegliere la chiave a brugola} dalla barra
        degli strumenti.
  \item Lo studente clicca sulla \textbf{vite} (\texttt{vite\_culatta\_1}).
  \item La vite si \textbf{svita} girando e uscendo di 30 centimetri.
  \item La vite si \textbf{appoggia a terra} in mezzo secondo.
  \item Lo studente clicca la \textbf{freccia avanti} per andare al prossimo step.
\end{enumerate}

\begin{figure}[H]
  \centering
  % \includegraphics[width=13cm]{immagini/cap20/sequenza_step_completo.png}
  \imgph{13cm}{Sequenza di 4 fotogrammi orizzontali dello step ``Rimuovi la prima vite'': 1) stato iniziale (vite avvitata, fumetto compare con highlight); 2) studente seleziona brugola, cursore cambia; 3) vite si svita (rotazione + traslazione); 4) vite appoggiata a terra accanto alla culatta. Numerazione 1--4 sotto.}
  \caption{Lo step ``Rimuovi la prima vite'' eseguito: cosa vede lo studente, momento per momento.}
  \label{fig:sequenza-step-completo}
\end{figure}

\section{Step automatici in catena}

\begin{esempiocvt}[Sequenza di viti consecutive]
\begin{lstlisting}
[step "Prima vite" highlight]
element     = vite_coperchio_1
tool        = hand
description = Seleziona il tool "Mano" e rimuovi la prima vite del coperchio.
do          : unscrew
do          : place duration 0.2

[step "Seconda vite" auto]
element     = vite_coperchio_2
description = Rimuovi la seconda vite del coperchio.
do          : unscrew
do          : place duration 0.2

[step "Terza vite" auto]
element     = vite_coperchio_3
description = Rimuovi la terza vite del coperchio.
do          : unscrew
do          : place duration 0.2
\end{lstlisting}
\end{esempiocvt}

\textbf{Cosa succede:}

\begin{itemize}
  \item La \textbf{prima} vite la deve fare lo studente (perché ha
        \val{highlight}, niente \val{auto}).
  \item Dopo che ha tolto la prima vite con successo, la \textbf{seconda} e la
        \textbf{terza} si svitano \textbf{da sole} (\val{auto}).
\end{itemize}

\begin{trucchetto}
\val{tool = hand} lo scrivi \textbf{solo nella prima vite}. Le successive con
\val{auto} lo \textbf{ereditano} dal contesto precedente.
\end{trucchetto}

\section{Esempio di drag \& drop}

\begin{esempiocvt}[Step con trascinamento]
\begin{lstlisting}
[step "Monta il coperchio" highlight]
description = Posiziona il coperchio sopra il filtro.
tool        = hand
drag        = coperchio distance 0.3
\end{lstlisting}
\end{esempiocvt}

\textbf{Cosa succede:}

\begin{enumerate}
  \item Apparirà il fumetto ``Posiziona il coperchio sopra il filtro''.
  \item Lo studente seleziona la mano.
  \item Lo studente \textbf{clicca e trascina} il coperchio.
  \item Quando il coperchio è entro 30 cm dalla sua sede originale,
        \textbf{si aggancia automaticamente} lì.
  \item Lo step viene completato e si va avanti.
\end{enumerate}

\section{Esempio di click su un pulsante 3D}

\begin{esempiocvt}[Step con click su pulsante]
\begin{lstlisting}
[step "Vai alla schermata MDI" highlight]
description = Premi il pulsante MDI sul pulpito.
element     = pulpito.Pulsante_mdi
after       : schermo = schermo002
\end{lstlisting}
\end{esempiocvt}

\textbf{Cosa succede:}

\begin{enumerate}
  \item Apparirà il fumetto ``Premi il pulsante MDI sul pulpito''.
  \item Il pulsante MDI si \textbf{illumina} di giallo (auto-evidenziato).
  \item Lo studente lo clicca.
  \item Lo schermo del pulpito \textbf{cambia} mostrando la schermata
        ``schermo002''.
  \item Si passa al prossimo step automaticamente.
\end{enumerate}

\section{Esempio di click su una porta}

\begin{esempiocvt}[Step con click su porta che si apre]
\begin{lstlisting}
[step "Apri la porta della cabina" highlight]
description = Apri la porta della cabina della macchina.
element     = a500.porta
do          : rotate around (0.615, 0, 0) by (0, 90, 0) duration 1.0
\end{lstlisting}
\end{esempiocvt}

\textbf{Cosa succede:}

\begin{enumerate}
  \item Apparirà il fumetto ``Apri la porta della cabina della macchina''.
  \item La porta si \textbf{illumina} di giallo per dire ``cliccami''.
  \item Lo studente la clicca.
  \item La porta \textbf{ruota} di 90 gradi attorno al cardine in 1 secondo.
  \item Quando l'animazione è finita, si passa al prossimo step.
\end{enumerate}

\begin{notablu}[Stessa keyword, comportamento diverso]
Nei due esempi sopra abbiamo la stessa parola \kw{element =}. Nel primo
caso (pulsante) lo step ha \kw{after :} e si passa subito al prossimo step
dopo il click. Nel secondo (porta) lo step ha \kw{do :} e si aspetta che
l'animazione finisca. \textbf{Tu non devi dire al programma quale dei due
casi è}: lui lo capisce dalle altre righe dello step.
\end{notablu}

%--------------------------------------------------------------------------
\chapter{Errori comuni e come evitarli}
\label{ch:errori}
%--------------------------------------------------------------------------

\begin{figure}[H]
  \centering
  % \includegraphics[width=13cm]{immagini/cap21/mappa_10_errori.png}
  \imgph{13cm}{Griglia 2$\times$5 con dieci celle, ognuna con un simbolo \sino{} rosso e una parola chiave breve: ``Maiuscolo'', ``Virgolette'', ``Numerare'', ``.glb'', ``Italiano'', ``Spazi'', ``\texttt{//}'', ``AutoSet'', ``ActiveButtons'', ``Animated...''. Sotto il titolo ``i 10 errori più comuni''.}
  \caption{La mappa rapida dei dieci errori più comuni: usa questa pagina come check-list quando un tutorial non funziona.}
  \label{fig:mappa-10-errori}
\end{figure}

\section{Errore 1: scrivere le righe con la lettera maiuscola}

\begin{sbagliato}[Sbagliato (vecchia sintassi v1/v2)]
\begin{lstlisting}
Element=vite_coperchio_1
Tool=Mani
\end{lstlisting}
\end{sbagliato}

\begin{giusto}[Giusto (sintassi nuova v3)]
\begin{lstlisting}
element = vite_coperchio_1
tool    = hand
\end{lstlisting}
\end{giusto}

Le maiuscole continuano a funzionare per \textbf{retrocompatibilità}, ma la
sintassi nuova vuole tutto \textbf{minuscolo}.

\section{Errore 2: dimenticare le virgolette nel nome dello step}

\begin{sbagliato}
\begin{lstlisting}
[step Rimuovi la prima vite]
\end{lstlisting}
\end{sbagliato}

\begin{giusto}
\begin{lstlisting}
[step "Rimuovi la prima vite"]
\end{lstlisting}
\end{giusto}

Senza virgolette il programma si confonde perché ci sono spazi nel nome.

\section{Errore 3: numerare le azioni manualmente}

\begin{sbagliato}[Non è necessario (vecchia sintassi)]
\begin{lstlisting}
Action1 = unscrew
Action2 = place(0.2)
Action3 = move by (0, 0, -0.1)
\end{lstlisting}
\end{sbagliato}

\begin{giusto}[Sintassi nuova]
\begin{lstlisting}
do : unscrew
do : place duration 0.2
do : move by (0, 0, -0.1) duration 0.5
\end{lstlisting}
\end{giusto}

Tu scrivi \kw{do :} e il programma le numera da solo. Se vuoi inserirne una
in mezzo, basta aggiungere una riga -- niente da rinumerare.

\section{Errore 4: scrivere il nome del file \texttt{.glb}}

\begin{sbagliato}[Non serve]
\begin{lstlisting}
element = models/vite.glb
\end{lstlisting}
\end{sbagliato}

\begin{giusto}
\begin{lstlisting}
element = vite
\end{lstlisting}
\end{giusto}

Il programma aggiunge da solo \texttt{models/} e \texttt{.glb}. Tu scrivi solo il
\textbf{nome del pezzo}.

\begin{notablu}[Eccezione: child interno a un modello]
Se vuoi un pezzo dentro un altro (es.\ la ``porta'' dentro il modello ``a500''),
si scrive con il \textbf{punto}:
\begin{lstlisting}
element = a500.porta
\end{lstlisting}
\end{notablu}

\section{Errore 5: usare i nomi italiani degli attrezzi}

\begin{sbagliato}[Vecchio stile]
\begin{lstlisting}
tool = Mani
tool = ChiaveBrugola
\end{lstlisting}
\end{sbagliato}

\begin{giusto}[Stile v3]
\begin{lstlisting}
tool = hand
tool = hex_key
\end{lstlisting}
\end{giusto}

I nomi italiani continuano a funzionare ma quelli inglesi sono più moderni.

\section{Errore 6: dimenticare gli spazi attorno all'\texttt{=}}

Il programma è abbastanza tollerante, ma per \textbf{leggibilità} scrivi sempre:

\begin{giusto}[Ottimo]
\begin{lstlisting}
element = vite_coperchio_1
tool    = hand
\end{lstlisting}
\end{giusto}

\begin{sbagliato}[Funziona ma è brutto]
\begin{lstlisting}
element=vite_coperchio_1
tool=hand
\end{lstlisting}
\end{sbagliato}

\section{Errore 7: mettere commenti con \texttt{//}}

\begin{sbagliato}[Da non fare -- vecchia sintassi deprecata]
\begin{lstlisting}
// Questo è un commento
\end{lstlisting}
\end{sbagliato}

\begin{giusto}[Stile v3]
\begin{lstlisting}
# Questo è un commento
\end{lstlisting}
\end{giusto}

Il \texttt{\#} è il modo moderno. Il \texttt{//} ancora funziona ma il programma
ti avvisa che è vecchio.

\section{Errore 8: usare \texttt{AutoSetVariant} invece di \texttt{after :}}

\begin{sbagliato}[Vecchio stile]
\begin{lstlisting}
[step "Scambio Utensile" machine]
AutoSetVariant=tool=tool0
\end{lstlisting}
\end{sbagliato}

\begin{giusto}[Stile v3]
\begin{lstlisting}
[step "Scambio Utensile" machine]
after : tool = tool0
\end{lstlisting}
\end{giusto}

\kw{after :} si occupa da solo di cambiare la variante: se lo step ha un
\kw{element =} cliccabile cambia dopo il click, altrimenti cambia quando lo
step parte.

\section{Errore 9: usare \texttt{ActiveButtons} + \texttt{AcceptTrigger} a mano}

\begin{sbagliato}[Vecchio stile esplicito]
\begin{lstlisting}
ActiveButtons=chiave0
AcceptTrigger_Physical=pulpito.chiave0
OnPhysicalTrigger=setVariant:chiave=chiave1;setVariant:schermo=schermo005
\end{lstlisting}
\end{sbagliato}

\begin{giusto}[Stile v3 -- 2 righe sole]
\begin{lstlisting}
element = pulpito.chiave0
after   : chiave = chiave1; schermo = schermo005
\end{lstlisting}
\end{giusto}

L'\kw{element} dice ``il pezzo da cliccare è la chiave'', e \kw{after :} dice
``quando lo studente la clicca, applica questi cambi di stato''. Le tre righe
\texttt{ActiveButtons}, \texttt{AcceptTrigger\_Physical} e
\texttt{OnPhysicalTrigger} le aggiunge automaticamente il programma.

\section{Errore 10: scrivere \texttt{button = ...} (vecchia versione del v3)}

\begin{sbagliato}[Vecchia versione v3 (deprecata)]
\begin{lstlisting}
button = pulpito.Pulsante_mdi
after  : schermo = schermo002
\end{lstlisting}
\end{sbagliato}

\begin{giusto}[v3 attuale -- una sola keyword]
\begin{lstlisting}
element = pulpito.Pulsante_mdi
after   : schermo = schermo002
\end{lstlisting}
\end{giusto}

Nelle prime versioni del v3 c'era \kw{button =} per i pulsanti 3D. Ora
\textbf{tutto è \kw{element}}: pulsanti, porte, leve, viti. Il programma
capisce da solo cosa fare guardando le altre righe dello step. La parola
\kw{button =} continua a funzionare per non rompere i tutorial vecchi, ma il
programma ti avvisa che è deprecata.

\section{Errore 11: usare \texttt{Animated...} invece di \texttt{animation}}

\begin{sbagliato}[Vecchio stile]
\begin{lstlisting}
ActiveButtons=Pulsante_r_unlock
AcceptTrigger_Physical=remote.Pulsante_r_unlock
AnimatedImagesFolder=screens/mandrino
AnimatedPosition=(60%,60%)
AnimatedMaxTriggers=8
AnimatedFrameDelay=20
\end{lstlisting}
\end{sbagliato}

\begin{giusto}[Stile v3 -- 2 righe sole]
\begin{lstlisting}
element   = remote.Pulsante_r_unlock repeat 8
animation = folder screens/mandrino at (60%, 60%) frame 20ms
\end{lstlisting}
\end{giusto}

\val{repeat 8} dice ``8 click totali''. \kw{animation =} dice cosa mostrare
nel frattempo.

%--------------------------------------------------------------------------
\chapter{Glossario delle parole importanti}
\label{ch:glossario}
%--------------------------------------------------------------------------

\begin{description}[leftmargin=2.5cm,style=nextline]

\item[\textbf{Block} (blocco)]
Un pezzo di file racchiuso tra \texttt{[} e \texttt{]}. Esempi:
\sez{[scene]}, \sez{[section "\ldots"]}, \sez{[step "\ldots"]},
\sez{[state \ldots]}, \sez{[object \ldots]}.

\item[\textbf{Section} (sezione, capitolo)]
Un grande pezzo di tutorial (es.\ ``Pulizia Filtro''). Contiene tanti step.

\item[\textbf{Step} (passo)]
Un singolo passo del tutorial. Lo studente vede il fumetto, fa l'azione,
va al prossimo.

\item[\textbf{Element} (elemento)]
Il pezzo 3D coinvolto nello step. \textbf{È sempre la stessa parola}, qualunque
sia il pezzo: una vite (\texttt{vite\_coperchio\_1}), un pulsante
(\texttt{pulpito.Pulsante\_mdi}), una porta (\texttt{a500.porta}), una leva. Il
programma capisce dal contesto cosa fare quando lo studente lo clicca.

\item[\textbf{Tool} (attrezzo)]
L'attrezzo che lo studente deve scegliere. Cinque possibili: \val{hand},
\val{hex\_key}, \val{wrench}, \val{air}, \val{spray}.

\item[\textbf{Description} (descrizione)]
Il testo che appare nel \textbf{fumetto} della spiegazione.

\item[\textbf{Action / \texttt{do :}}]
Un'azione che il pezzo fa nello step. Possono essercene tante in fila.

\item[\textbf{Camera}]
Gli ``occhi'' dello studente. Si dirige con
\val{camera = position (\ldots), target (\ldots), \ldots}.

\item[\textbf{Drag}]
Lo studente trascina il pezzo con il mouse fino alla destinazione.

\item[\textbf{After}]
Cosa succede subito dopo che lo studente clicca l'\kw{element} (cambia uno o
più stati con \texttt{after : nome\_state = variante}).

\item[\textbf{Hold}]
Lo studente prende un oggetto in mano (es.\ un telecomando).

\item[\textbf{Bandierine}]
Parole speciali alla fine del nome di uno step. Tre esistono:
\val{auto}, \val{machine}, \val{highlight}.

\item[\textbf{Coordinate (x, y, z)}]
Tre numeri che dicono ``dove'' nello spazio. X = destra/sinistra, Y = su/giù,
Z = avanti/indietro.

\item[\textbf{Origin} (suffisso \texttt{\_original})]
Un riferimento alla \textbf{posizione iniziale} di un pezzo. Esempio:
\texttt{tappino\_grasso\_dx\_original} significa ``la posizione dove sta il
tappino di destra all'avvio''.

\item[\textbf{Companion} (compagno)]
Un pezzo che si muove \textbf{insieme} a quello principale dello step.
Esempio: quando svitate la flangia, anche il tubo del grasso la segue.

\item[\textbf{Variant} (variante)]
Una ``versione'' diversa di un pezzo. Esempio: lo schermo del pulpito può
essere \texttt{schermo001} (home) o \texttt{schermo002} (MDI). Cambiare
variante significa mostrare una versione diversa.

\item[\textbf{State} (stato)]
Un gruppo di varianti possibili. Esempio: lo state \texttt{schermo} può
essere uno di [\texttt{schermo001}, \texttt{schermo002}, \ldots].
Si dichiara nel blocco \sez{[state nome]} dentro \file{objects.ini}.

\item[\textbf{Object} (oggetto interattivo)]
Un modello 3D con pulsanti e parti che lo studente può cliccare
(es.\ il pulpito). Si dichiara nel blocco \sez{[object nome]} dentro
\file{objects.ini}.

\item[\textbf{Position} / \textbf{Rotation}]
Righe che spostano o ruotano un pezzo \emph{prima} che lo step parta. Si
scrivono in \sez{[scene]}, in una \sez{[section]} o dentro uno \sez{[step]}.

\item[\textbf{Companion}]
Un pezzo che si muove \textbf{insieme} a quello principale dello step. Può
seguirlo passivamente (\val{follows master}) o fare il proprio movimento
(\val{moves by}).

\item[\textbf{Repeat N} (in \texttt{element = X repeat N})]
Lo studente deve cliccare l'\kw{element} N volte. Si usa con \kw{animation =}
per animazioni a frame (apri/chiudi la pinza, sfoglia pagine, ecc.).

\item[\textbf{Animation}]
Una sequenza di immagini 2D mostrate sopra la scena 3D. Avanza ad ogni clic
del pulsante (vedi \val{repeat}).

\item[\textbf{View} (vista di un display)]
La pagina mostrata da uno schermo digitale 3D. Si cambia con \kw{view = nome}.
Le pagine possibili sono dichiarate in \file{objects.ini}.

\item[\textbf{Hold}]
Lo studente prende un oggetto in mano (\val{hold = pick}), continua a tenerlo
(\val{hold = held}) o lo lascia (\val{hold = release}).

\item[\textbf{Message / Title / Video / Image}]
Mostra una finestra modale con testo, titolo, video o immagine. Lo studente
chiude con OK e va avanti.

\end{description}

%--------------------------------------------------------------------------
\chapter*{Conclusione}
\addcontentsline{toc}{chapter}{Conclusione}
%--------------------------------------------------------------------------

Adesso sai scrivere un tutorial. Ricorda i \textbf{tre segreti}:

\begin{enumerate}
  \item \textbf{Scrivi minuscolo} (per la sintassi v3).
  \item \textbf{Le virgolette} nei nomi di section e step.
  \item \textbf{\kw{do :}} per ogni azione, una per riga, in ordine.
\end{enumerate}

Se ti dimentichi qualcosa, riapri questo manuale al \textbf{glossario}
(Capitolo~\ref{ch:glossario}) o all'\textbf{esempio completo}
(Capitolo~\ref{ch:esempio}).

\vfill
\begin{center}
\textit{\Large Buon divertimento e buon tutorial!}
\end{center}

\end{document}